% 本文件由 LaTeX 排版 Agent 根据有效 translation.json 自动生成。
% author=Tertullian; work_id=0316; title=论肉身的复活

\chapter{论肉身的复活}

% structure: p0001 source=h1 target=omitted reason=page-title-already-rendered-by-parent

\EditorialNote{本书所针对的异端者，正是那些主张德谬哥——即创造这个世界并赐予梅瑟律法的神——与至高的天主相对立的人。因此，他们认为德谬哥的一切工作都带有内在的败坏与无价值，其中也包括人的肉身或身体；他们断言肉身不能复活，唯独灵魂才能承继不朽。}\par

% structure: p0003 source=h2 target=section reason=relative-heading-rank; base=h2

\section{第一章　肉身复活之道理由福音显明。异教中偶尔有类似之微弱闪光。异教教导的矛盾。}

死者的复活是基督徒的信靠。因着它，我们成为信徒。真理迫使我们相信这一信德道理——这真理是天主所启示，却被群众嗤笑，他们以为死后一无所存。然而，他们仍向死者致敬，且\Emph{那样}也是按照遗嘱以最昂贵的方式，并用季节所能供应的最精美宴席来款待他们，这是基于一种假设：那些被他们宣称毫无知觉的人仍然保有食欲。但（任凭群众嗤笑吧）：我反而更要嗤笑他们，尤其当他们以最残忍的不人道焚化死者，随后又立即以贪饕的饱足来纵容他们，用同一火焰既尊崇又侮辱他们。这是何等虔诚，竟以残忍戏弄\Emph{其牺牲}？群众焚烧祭品给那些已被焚烧者，这是献祭还是侮辱？然而，智者有时也与平民百姓同流合污。依照伊比鸠鲁学派，死后一切皆无。塞内卡也说，死后万物终结，连死亡本身也归于无有。令人欣慰的是，毕达哥拉斯、恩培多克勒以及柏拉图主义者的哲学也同样重要，持相反观点，宣称灵魂不朽；并且以最接近我们教义的方式断言，灵魂确实返回肉身，尽管不是同一个肉身，甚至不总是人的肉身：因此，欧佛尔波被认为转化成了毕达哥拉斯，荷马则成了一只孔雀。他们坚定地宣告灵魂在肉身中的更新，认为改变（肉身状态的）性质比完全否认它更可容忍：他们至少叩响了真理的门，却没有进入。这样，世界虽充满错误，却并非不知死者的复活。\par

% structure: p0005 source=h2 target=section reason=relative-heading-rank; base=h2

\section{第二章　犹太人的撒杜塞人是异教哲学家与异端者之间关于此道理的连接。此道理的根本重要性被肯定。在异端者的估计中，灵魂的未来状态比肉身更受优待。然而，灵魂的灭绝曾被一位名叫卢坎的主张。}

既然在天主的\Emph{教会}内，甚至有一个派别与伊壁鸠鲁派比与先知们更为接近，这就给我们机会知道基督对（所说派别，即）撒杜塞人如何评价。因为揭露先前隐藏的一切，赋予可疑之点以确定性，成就人只有预尝之事，使预言的对象成为现世现实，\Emph{并且}不仅藉祂自己，而且在祂自己内提供死者复活的确实证据，这些事都是保留给基督的。然而，我们现在必须准备应对的是另一种撒杜塞人，他们仍分沾其教义。例如，他们承认复活的半数；即仅仅灵魂的复活，鄙视肉身，正如他们也鄙视肉身的主自己。确实，没有别人拒绝承认身体的实体从死亡复原，只有异端者——第二神祇的发明者。因此，他们被迫给基督一种不同的安排，以免祂被认为属于创造主，他们便在祂肉身的\Emph{方面}取得了第一个错误；与\Person{玛尔西翁}和\Person{巴西里德}一致争辩说，肉身没有实体；或按照\Person{瓦伦提努斯}的异端学说和\Person{阿佩莱斯}，认为肉身具有其特有的品质。因此，他们就把基督没有分有的那个实体排除在一切从死亡复原之外，自信地认为这构成了反对复活的最强依据，因为肉身已经在基督内复活了。因此，我们先前出版了我们的著作\WorkTitle{论基督的肉身}；在其中，我们既提供祂肉身实体的证据，反对认为它是虚影的意见；也为它主张一种毫无特殊状况的人性——正是这样的本性使基督成为人且为人子。因为当我们证明祂穿戴了肉身并处于身体状态时，我们同时就驳斥了异端，确立了这样的准则：必须相信除了创造主之外没有别的存有是天主，因为我们表明基督——在其中清楚看见天主——恰恰是创造主所应许的那样本性。他们在天主是创造主、基督是肉身的救赎者这两点上如此被驳倒后，他们立刻也将在肉身复活上失败。确实，没有比这更合理的方法。我们断言，与异端者的争论在多数情况下应按此方式进行。因为适当的方法要求，结论应总是从最重要的前提出发，以便就关键点先达成一致，藉此可以认为特定的问题已被决定。因此，异端者由于自觉软弱，从不按有序方式进行讨论。他们很清楚，要暗示存在第二位神抵毁世界的创造主是多么困难的任务——那位创造主是万民由祂工程的见证自然认识的，祂在\Emph{祂的存在}的奥秘中超越万有，并特别在先知中彰显；然后，他们以考虑更紧迫的探究——即人自身的救恩——为借口，开始质疑复活；因为相信肉身复活比相信天主独一更困难。这样，他们在剥夺了讨论逻辑顺序的优势，并用贬低肉身的可疑暗示使之难堪之后，逐渐将论点引向接受第二位神，从而摧毁并改变了我们希望的基础。因为一旦人从他对创造主所抱的确定希望中堕落或远离，他就容易被引向另一个希望的对象，然而凭其自身，他很难不怀疑这个对象。其实，正是通过应许间的差异，暗示了神祇的不同。我们看见多少人这样被网罗，他们在肉身复活上被击败，而未能坚持天主独一的论点！因此，对于异端者，我们已经表明应当用什么武器迎战他们。并且事实上，我们在分别针对他们的论文中已经与他们交锋：在反对\Person{玛尔西翁}的著作中论及惟一天主与祂的基督，在反对四种异端的书中论及主的肉身，都是为了特别开启当前探究的道路；因此我们现在只需讨论肉身复活，就好像它对我们自己也是不确定的，即在创造主的安排中。因为许多人未受教育；更多人信心摇动，还有一些人意志薄弱：这些人必须被教导、指引、加强，因为天主的独一性与我们教义的维护将一并被辩护。因为如果肉身复活被否认，\Emph{那个首要的信德条款}就被动摇；如果被确认，它就被确立。我想，没有必要讨论灵魂的得救；因为几乎所有的异端者，无论他们如何构想它，当然都不否认\Emph{那}。我们可以忽略某个\Person{路卡努斯}，他甚至不放过我们本性的这一部分，他追随\Person{亚里士多德}将其归结为消散，并以某种别的东西取代之。按照他的说法，是某种第三本性要复活，既非灵魂也非肉身；换言之，不是人，而可能是熊——例如，\Person{路卡努斯}本人。连他也在我们关于灵魂全部状态的书中受到了我们充分的关注，我们在那里维护灵魂的特殊不朽，同时既承认唯独肉身的消散，又强调主张肉身的复原。在那部著作的主体中，收集了我们因偶然原因的压力而在别处不得不保留的一切要点。因为正如我的习惯，在问题首次出现时只轻微触及，所以我也必须推迟对它们的考虑，直到轮廓可以用完整的细节填充，并且被推迟的问题按其自身价值被处理。\par

% structure: p0007 source=h2 target=section reason=relative-heading-rank; base=h2

\section{第三章  外邦人也持有某些真理。然而，他们常在宗教信仰和道德实践上犯错。不可效法外邦人对基督奥迹的无知。异端者却反常地倾向效法他们。}

诚然，人凭其本性在有关天主的事上也能有智慧，但只为见证真理，而非维护谬误；只当他依照天主眷顾而行、而非违背之时。因为有些事连本性也能知晓：例如许多人相信灵魂不死；我们天主的认识是人人皆有的。因此，我可以使用\Person{柏拉图}的意见，当他宣称：\Quote{每一个灵魂都是不死的。}我也可以使用一个民族的良心，当它见证万神之神。同样，我可以使用我们共同本性的其他理智，当它们宣告天主是审判者。他们说：\Quote{天主都看见}；并且\Quote{我把你交托给天主}。但当他们说：\Quote{已死之物便是死了}；\Quote{活着时就享受人生}；\Quote{死后万有皆终，连死亡本身也不例外}，那时我必须记住：按天主的判断，\Quote{人的心是灰烬}\BibleRef{依 44:20}，而且\Quote{这世界的智慧}本身也被（受默感的话）称为\OriginalTerm{愚妄}{foolishness}。然后，如果连异端者也在凡俗堕落的思绪或世界的想象中寻求避难所，我必须对他说：异端者啊，与外邦人分道扬镳吧！因为虽然你们在想象一位天主上全都一致，但既然你们以基督之名这样做，只要你们自以为基督徒，你们就与外邦人不同：把外邦人自己的看法还给他，因为他并不从你那里学习。如果你自己有眼睛，为何倚靠盲目的向导？如果你已穿上基督，为何让赤身者给你披衣？既然宗徒赐给你自己的盔甲，为何使用别人的盾牌？他不如从你那里学会承认肉身的复活，而不是你从他那里学会否认它；因为如果基督徒必须否认它，他们凭自己的知识就足够了，无需从无知的群众那里得到教导。因此，否认基督徒所承认的这一教义的人，将不会是基督徒；更何况他否认的依据是一个非基督徒所采纳的理由。的确，将异端者与外邦人共享的智慧夺去，让他们只凭圣经支持他们的探询：他们便无法立足。因为使人常识可取的正是它的单纯、它参与同样的感受、以及它意见的共通性；它越是被认为可信，因为其断论赤裸、公开、众所周知。相反，天主的理致存在于事物的精髓和骨髓之中，不在表面，并且常常与表象相悖。\par

% structure: p0009 source=h2 target=section reason=relative-heading-rank; base=h2

\section{第四章  外邦人与异端者同样诋毁肉身及其功能：对最终恢复如此软弱卑微之质料常有的非难。}

因此，异端者立即从这一点出发，描绘他们教义的初稿，而后添上细节，深知人的心思多么容易被其影响所掳，（且受）那有利于他们设计的人类共同情感所驱动。从异端者那里，你能听到什么比从外邦人那里更早、更广泛的言论？难道他们的（重担）从起初就不仅是对肉身的责难——针对它的起源、它的质料、它遭遇的不幸和无可避免的结局——从最初用地上的泥土成形就是不洁的，以后因它种子传递的泥泞而更不洁；无价值、软弱、满覆罪债、充满悲惨、多受苦楚；而在这一切贬抑的记录之后，它归于原来的土地和尸首之名，并且注定甚至从这个可憎的名称逐渐消失，直到此后一切名称归于虚无？你无疑是个精明人：那么，你能说服自己，这肉身一旦从视线、触觉和记忆中消失，能再从不腐坏恢复到完整，从破碎状态到坚实状态，从空虚状态到丰满状态，从虚无到某物——让吞灭的火、海水、野兽的胃、飞鸟的嗉囊、鱼的肚子以及时间本身的巨大胃口都将它交还？那已经朽坏的肉身，会如此被期望复原，以致瘸子、独眼、瞎眼、癞病、瘫痪者都将回来，尽管他们回到旧状态毫无乐趣？或者他们会痊愈，因此不得不害怕遭受那样的苦难？在那种情况下，（我们必须说）肉身的恢复有什么后果？它会不会再次受制于它现在的一切需要，尤其是食物和饮料？我们是否要用肺漂浮（在空中或水中），肠中疼痛，用羞耻的器官并不羞愧，用所有肢体劳苦？难道必须再有溃疡、伤口、发烧、痛风，并且再一次渴望死亡？当然，这些将是肉身恢复中伴随的渴望，不过是重复想脱离它的欲望。好了，这一切我们都以非常柔和含蓄的措辞述说了，以配合我们文风的特性；（但要知道）这些人实际使用何等不雅的言辞，你必须在他们的会议中考验他们，无论是外邦人还是异端者。\par

% structure: p0011 source=h2 target=section reason=relative-heading-rank; base=h2

\section{第五章  赞美肉身的几点回应。肉身由天主创造。人的身体实际上先于他的灵魂。}

因此，既然所有未受教育的人仍然根据这些常识性的观点来形成他们的意见，而那些动摇不定和意志薄弱的人又因这些相同的观点而重新陷入困惑；并且，那指向我们的第一攻城槌就是粉碎肉身状况的观点，因此，我们必须相应地组织我们的防御，以便首先保卫肉身的状况，通过我们的颂赞来击退他们对肉身的贬低。因此，异端者们要求我们不仅运用哲学，也同样运用修辞学来对付他们。那么，关于这个脆弱、贫穷、毫无价值的身体——他们甚至毫不犹豫地称其为邪恶，即使它是天使的作品，正如\Person{梅南德}和\Person{马库斯}所乐意认为的，或者是如\Person{阿佩莱斯}所教导的，某个火性的存在、一位天使的塑造——仅仅因为身体得到了一个次级神祇的支持和保护，就足以确保身体受到尊重。我们知道，天使的等级仅次于\Emph{天主}。现在，不管每个异端者所认为的至高天主是谁，我都可以公平地从那位意欲产生肉身的祂那里推导出肉身的尊严。因为，当然，如果祂不希望它产生，当祂知道它正在进行时，祂就会阻止它。因此，即使按照他们的原则，肉身同样也是天主的工作。没有作品不属于那允许它存在的祂。事实上，这是一个幸运的情况，他们的大部分教义，甚至包括最严苛的，都将整个人类的形成归于我们的天主。祂是多么伟大，你们这些相信祂是唯一天主的人非常清楚。那么，让肉身开始给你带来喜悦吧，因为它的创造者是如此伟大。但是，你说，连世界都是天主的工作，然而\BibleQuote{这世界的格式正在逝去，}\BibleRef{格前 7:31}正如宗徒自己所见证的；也不能因为世界是天主的工作就预先断定它会得到恢复。而且，如果宇宙在其毁灭之后不再被重新塑造，那么它的一部分为何要被重塑呢？如果你认为部分与整体平等，那你是对的。但我们坚持认为存在差异。首先，因为万物都是借着天主的圣言造的，没有祂，什么也没有造。 \BibleRef{若 1:3}现在，肉身也是从天主的圣言获得存在的，因为原则是，没有圣言就没有任何东西存在。\BibleQuote{让我们造人，}\BibleRef{创 1:26}祂在创造人之前说，并补充说，\Quote{用我们的手，}为了人的卓越性，以便他不与其他受造物相比。而\Quote{天主}，（圣经说）\BibleQuote{创造了人。}\BibleRef{创 1:27}毫无疑问，在过程中存在巨大的差异，这当然源于事物的本质。因为被造的受造物比它们为之而造的人低等；它们是为人类造的，后来被天主置于人的权下。因此，那些注定要服从的受造物，凭着\Emph{神圣}声音的命令、指示和独一的能力而出现，这是正确的；而人，相反，由于注定要成为它们的主，是由天主亲自塑造的，目的是使他能够行使他的主权，因为他是被主\Emph{自己}创造的。记住，人也被恰当地称为\Emph{肉身}，这个名称在人的命名中优先占有：\Quote{天主用地上的泥土造了人。}他现在成了人，而之前是泥土。\BibleQuote{祂将生命的气息吹在他的脸上，人（即泥土）就成了一个有灵的生命；天主将祂所造的人安置在园子里。}\BibleRef{创 2:7}因此，人起初是泥土，之后才成为完整的\Emph{人}。我希望你注意这一点，以便你知道，无论天主向人计划或应许了什么，都不仅仅是由于灵魂，也是由于肉身；即使不是出于他们起源的共性，也至少是由于后者在其名称中所拥有的\Emph{特权}。\par

% structure: p0013 source=h2 target=section reason=relative-heading-rank; base=h2

\section{第六章。在衡量肉身的卓越时，必须牢记的，不是材料的卑微，而是造物主的尊贵与技巧。基督分享了我们的肉身。}

让我们继续探讨眼前的问题——我若能成功地为肉身辩护，如同天主所赋予它的一切荣耀，它那时便已自夸，因为那微不足道、可怜的黏土落入了天主的手中——无论那手是什么——只被触摸一下便已足够幸福。但这\Emph{夸耀}为何？难道黏土在天主触摸之下，无需进一步劳作便立即成形？事实是，有一件大事正在进行，我们所思虑的受造物正被塑造出来。因此，它每次经历天主的手，被触摸、被拉扯、被拉伸、被塑形，便如此频繁地接受荣耀。想象天主完全沉浸其中——祂的手、祂的眼、祂的劳作、祂的旨意、祂的智慧、祂的眷顾，尤其是祂的爱，正规划着（这受造物）的轮廓。因为，无论当时创造者给予黏土何种形态和表情，基督已在祂的思念中，有朝一日将成为人，因为圣言也要成为黏土和肉身，正如当时的土。圣父曾对圣子这样说：\BibleQuote{让我们按我们的肖像，按我们的模样造人。}\BibleRef{创 1:26} 天主造了人，即祂所捏塑、所塑造的受造物；是按天主的肖像——换言之，是按基督的肖像——造了他。圣言也是天主，祂虽具有天主的肖像，却\BibleQuote{不以自己与天主同等为强夺的。}\BibleRef{斐 2:6} 因此，那当时就已披上基督——那将要降生在肉身中的基督——肖像的黏土，不仅是天主的工程，也是天主的抵押和保证。把\Emph{土}这个名称当作卑贱低劣的元素传来传去，意图玷污肉身的起源，又有什么意义呢？即使有其他任何材料可用于造人，也应当考虑造物主的尊荣，祂甚至通过选择材料而认为它值得，并通过祂的处置使之配得。菲狄亚斯之手用象牙塑造了奥林匹斯的朱庇特；人们崇拜\Emph{那雕像}，不再将它视为由最愚蠢的动物所成的\Emph{神}，而是世界至高之神——不是因为大象的庞大体型，而是因为菲狄亚斯的名声。难道永生的天主、真天主，不能通过祂自己的作为，把祂材料的任何卑劣之处清除干净，并治愈它的一切软弱吗？还是必须留下这一点，\Emph{以显示}人制造神比天主造人更高贵？如今，虽然黏土令人反感（因其卑劣），但它已变成别的东西。我拥有的是肉身，不是土，即使肉身曾被说：\BibleQuote{你是尘土，仍要归于尘土。} 这些话提及的是起源，而非重提质料。这特权已赐给\Emph{肉身}，使其比其起源更高贵，并通过其转变而获得更大的幸福。如今，甚至金子也是土，因来自土；但变成金子后，它便不再是土，而是一种远为不同的质料，更璀璨、更高贵，尽管源自相对黯淡、卑微的来源。同样，天主完全可以清除我们肉身金子的所有污点——你们认为是来自其\Emph{原生}黏土的污点——通过炼净原始材料的渣滓。\par

% structure: p0015 source=h2 target=section reason=relative-heading-rank; base=h2

\section{第七章 创造肉身的土质材料，经天主的塑造而奇妙提升；因灵魂加入人的构造，它成为创造中的首要工程。}

但或许肉身的尊贵似乎被贬低，因为它并未像泥土\Emph{起初}那样被天主亲手操控。然而，当天主为着从泥土中生出肉身的目的而特意处理泥土时，祂是为此肉身而操劳的。此外，我还要你们知道，肉身是在何时、以何种方式从\Emph{其}泥土中绽放出美丽。因为，正如某些人所认为的那样，\Person{亚当}和\Person{厄娃}在失去乐园时所穿的\Quote{皮衣}\BibleRef{创 3:31}本身正是泥土形成肉身的过程，但早在之前，亚当已认出女人身上的肉身是他自己本体的繁衍（\Quote{这才是我骨中的骨，肉中的肉}\BibleRef{创 2:23}），而且从男人身上取出女人之事是以肉身完成的；但我认为，如果亚当仍是泥土，那么它本应以泥土来补充。因此，泥土被消除并吸收为肉身。\Emph{这是何时发生的？}就在人因天主的嘘气成为有灵的生物之时——正是那能将泥土硬化成另一种质料（如陶器）的嘘气，如今则成为肉身。同样，陶匠也有能力通过调节火焰的吹息，将陶土材料改变为更坚硬的材料，并塑造出比原质料更美的形式，且拥有其自身的种类和名称。因为，尽管圣经说：\Quote{泥土岂能对陶匠说？}\BibleRef{罗 9:20}即人岂能与天主\Emph{争论}？虽然宗徒提到\Quote{瓦器}\BibleRef{格后 6:7}，他指的是人，原本是泥土。而这器皿就是肉身，因为肉身是由天主的嘘气从泥土中\Emph{造成的}；后来它被\Quote{皮衣}覆盖，即覆盖在它上面的皮肤外衣。事实正是如此：如果你剥去皮肤，就露出肉身。因此，被剥下成为战利品的东西，当它仍覆盖着时就是一件外衣。因此，宗徒称割礼为\Quote{脱去肉身}\BibleRef{哥 2:11}，就肯定了皮肤是外衣。既然如此，你既已有被天主之手荣耀的泥土，又有因天主嘘气而更荣耀的肉身，凭借这嘘气，肉身不仅抛弃了泥土的原质，还披上了灵魂的装饰。\Emph{你}难道比\Emph{天主}更谨慎吗？你确实不愿将\Place{斯基泰}、\Place{印度}的宝石和\Place{红海}的珍珠镶嵌在铅、铜、铁甚至银中，而是将它们镶嵌在最为珍贵和精工的金器中；或者，你又为最上等的葡萄酒和最昂贵的香膏预备最相配的器皿；同样，你为精炼的剑配以同等价值的剑鞘；而天主却要将祂自己灵魂的影子、祂自己圣神的嘘气、祂自己口中的作为，交付给某种最卑贱的鞘，以此卑劣的交付当然注定其受谴责。那么，祂是否已将灵魂置于肉身中，甚或插入并混合其中？是的，而且结合如此紧密，以至于难以确定是肉身承载灵魂，还是灵魂承载肉身；或是肉身作为灵魂的\Emph{侍从}，还是灵魂作为肉身的\Emph{侍从}。然而，更可信的是灵魂被服务并拥有主权，因其本性更接近天主。\BibleRef{若 4:24}这甚至给肉身增添了荣耀，因为它既包含了最接近天主的本质，又使自己分享了灵魂实际的主权。因为，有什么自然的享受、世界的产物、元素的滋味，不是借着身体传递给灵魂的呢？否则怎么可能？灵魂不正是借助身体由全部感官（视觉、听觉、味觉、嗅觉、触觉）来支持的吗？灵魂不正是借助身体被洒上了神圣的能力，因为没有什么不是通过其言语能力实现的，即便是默示时也是如此？而言语是肉身器官的结果。技艺来自肉身；通过肉身，心灵的追求和能力也得以实现；一切工作、事业和生活的职责都由肉身完成；灵魂的活行动完全是肉身的工作，以至于灵魂停止活行动就等于与肉身分离。同样，死亡本身也是肉身的功能，正如生命的过程一样。现在，如果万物都通过肉身从属于灵魂，那么它们同样也属于肉身。作为你享受的媒介和工具，它也必须成为你享受的参与者和分享者。因此，被视为灵魂的仆役和侍从的肉身，结果也成为它的伴侣和同承继者。如果在暂时的事物中如此，为何在永恒的事物中不也如此呢？\par

% structure: p0017 source=h2 target=section reason=relative-heading-rank; base=h2

\section{基督教对肉身的照顾赋予它最大的荣耀。我们宗教的特权与我们的肉身最为密切。肉身也在宗教的职责和祭献中占有很大份额。}

现在，我愿就人之本性的普遍状况，提出这些为肉身辩护的评论。让我们思考它与基督信仰的特殊关系，看看天主赐予这卑微无用的物质多大的特权。诚然，只需说：除非灵魂在肉身内相信，否则它绝不能获得救恩；肉身乃是救恩所系的条件，此话千真万确。既然灵魂因救恩而被拣选事奉天主，实际上是肉身使灵魂能够如此事奉。的确，肉身被洗涤，好使灵魂得以洁净；肉身被傅油，好使灵魂得以祝圣；肉身被（十字架）标记，好使灵魂也得到坚固；肉身因覆手而成阴影，好使灵魂也被圣神光照；肉身领受基督的体血，好使灵魂也因\Emph{自己的}天主而得饱饫。既然它们在事奉中结合，在赏报上便不能分离。此外，那些中悦天主的祭献——我是指灵魂的战斗、斋戒、克己以及附属于这种义务的谦卑——是肉身一再为它自身的特殊痛苦而做的。同样，童贞、寡居、在隐密中对婚姻之床的节制，以及唯一的婚姻采纳，都是肉身善行所献给天主的馨香祭品。来吧，告诉我，当肉身为基督的名争战时，被拖到众人面前，暴露在所有人的仇恨中；当它在监狱中因残酷的光照剥夺、与世隔绝、在污秽、肮脏和难闻的食物中消瘦，甚至在睡眠中也不得自由，因为它被捆绑在铺上，在草褥上被折磨；当它最终在众目睽睽之下，受尽各种能想出的酷刑，并在痛苦中耗尽，为基督而死，献上最后的生命——多次在祂自己的十字架上，更不用说还有更残忍的酷刑手段——你对这肉身作何感想？最蒙福、最荣耀的，必定是那能如此完全地酬报其主人基督如此巨大债务的肉身，以至于唯一尚欠祂的，就是它应\Emph{因死亡}而不再欠祂更多——\Emph{即使在感恩中}也更有亏欠，因为它（永远）得了释放。\par

% structure: p0019 source=h2 target=section reason=relative-heading-rank; base=h2

\section{第九章 天主对人的肉身之爱，在基督对其恩宠中彰显。肉身是显示天主的慷慨与大能的最佳媒介。}

那么，再总结一下：那由神圣造物主亲手按天主的肖像所塑造的肉身；那被祂以自己的\Emph{生气}、按祂自己生命活力的样式所赋予生命的肉身；那被祂置于祂一切手造之物之上，使之居住、享受并统治它们的肉身；那被祂以圣事和训诲所装饰的肉身；祂喜爱其纯洁，赞许其克苦；祂视其为自己所受之苦为宝贵的肉身——（我说，这肉身）如此多次被带到天主面前，难道不会复活吗？断乎不可，断乎不可，（我重复），祂岂可把自己亲手的工作、自己思想的关切、自己圣神的居所、受造界的皇后、自己慷慨恩赐的承继者、自己宗教的女司祭、自己见证的斗士、自己基督的姊妹，弃置于永久的毁灭！我们以经验认识天主的良善；从祂的基督我们得知祂是独一的天主，而且是至善的。既然祂要求我们在爱祂之后爱近人如己，\BibleRef{玛 22:37}那么祂自己也会行祂所命令的。祂将爱那如此紧密且以多种方式作祂近人的肉身——（祂将爱它），虽然软弱，因为祂的能力在软弱中得以全备；\BibleRef{格后 12:9}虽然紊乱，因为\Quote{健康的人不需要医生，而是有病的人需要;} \BibleRef{路 5:31}虽然不体面，因为\Quote{我们对于不体面的肢体，越发放上相称的体面;} \BibleRef{格前 12:23}虽然失丧，因为祂说：\Quote{我来是为寻找并拯救失丧的人;} \BibleRef{路 19:10}虽然有罪，因为祂说：\Quote{我宁愿罪人得救，而不愿他死亡;} \BibleRef{则 18:23}虽然被定罪，因为祂说：\Quote{我击伤，也必医治。} \BibleRef{申 32:39}为何要指责那些等候天主、仰望天主、从天主得尊荣、蒙天主扶助的肉身所遭遇的种种情形呢？我敢说，如果这些灾祸从未临到肉身，那麼天主的慷慨、恩宠、仁慈（以及）一切有益的权能，就没有机会施展了。\par

% structure: p0021 source=h2 target=section reason=relative-heading-rank; base=h2

\section{第十章 圣经颂扬肉身，论其本质与前瞻。}

你持守那些贬低血肉的圣经；也接受那些高举血肉的圣经。你读任何贬抑它的段落；也请将你的目光转向那提升它的段落。\BibleQuote{一切血肉都是草。} \BibleRef{依 40:7} 然而，依撒意亚并不满足于只说这一点；他也宣告说：\BibleQuote{一切血肉都要看见天主的救恩。} \BibleRef{依 40:5} 他们留意天主在创世纪中说：\Quote{我的圣神不永远留在这些人中，因为他们是血肉；} 但随后又听见天主藉岳厄尔说：\BibleQuote{我要将我的圣神倾注在一切血肉身上。} \BibleRef{岳 3:1} 就连宗徒也不应该只凭他惯常指责血肉的一句话而被认识。因为尽管他说：\BibleQuote{在他肉体内没有善；} \BibleRef{罗 8:18} 尽管他肯定：\BibleQuote{凡随从血肉的人都不能中悦天主，} \BibleRef{罗 8:8} 因为\BibleQuote{血肉与圣神相争；} \BibleRef{迦 5:17} 然而在他所作的这些及类似断言中，受责备的不是血肉的\Emph{本体}，而是它的\Emph{行为}。此外，我们将在别处有机会指出，对血肉的任何指责若不同时导致对灵魂的谴责，就不能公平地加于血肉，因为灵魂迫使血肉执行其命令。不过，让我在此补充一点：在同处，保禄\BibleQuote{在身上带着主耶稣的印记；} \BibleRef{迦 6:17} 他也禁止我们的身体被亵渎，因为是\BibleQuote{天主的殿；} \BibleRef{格前 3:16} 他使我们的身体成为\BibleQuote{基督的肢体；} \BibleRef{格前 6:15} 并且劝勉我们高举并\Quote{在身体上光荣天主。} 因此，如果血肉的卑微阻碍了它的复活，为何它崇高的特权不能反而促成复活呢？——因为更符合天主的性格的是：将祂一度拒绝的恢复至救恩，而不是将祂一度认可的交予灭亡。\par

% structure: p0023 source=h2 target=section reason=relative-heading-rank; base=h2

\section{第十一章 天主的能力足以实现血肉的复活}

至此我已论及对血肉的赞颂，以反驳其仇敌——然而这些仇敌同时亦是其最大的朋友；因为没有谁像那些否认血肉复活的人那样生活在血肉中，因为他们轻视血肉的一切纪律，同时不相信血肉的惩罚。护慰者藉女先知\Person{普黎斯卡}的口所说的关于这些人的话是一句尖锐的箴言：\Quote{他们是属血肉的，却憎恨血肉。} 那么，既然血肉拥有了确保救恩赏报的最大保证，我们岂不应该也仔细思量天主本身的大能、威力与能力吗？——看祂是否如此伟大，足以重建并恢复那已败坏、堵塞且处处脱节的血肉大厦？——祂是否在自然的公开领域中颁布了任何类比，以使我们信服祂在这方面的能力，免得有人仍渴求认识天主，而对祂的信德除了相信祂全能之外别无基础？毫无疑问，在你们的哲学家中有人主张世界没有起始也没有创造者。然而，几乎所有异端都承认世界有起源和创造者，并将其创造归于我们的天主，这却是更为真实的事。因此，要坚信祂从无中创造了万物，这样你便因相信祂拥有如此大能而获得了对天主的认识。但有些人起初因他们对原质的看法而过于软弱，不能相信这一切。他们宁愿像哲学家那样，认为宇宙起初是天主从某种原质中创造的。现在，即使这种观点可以被真实地持有，既然必须承认天主在改造原质时产生了与原质本身所具有的截然不同的实体和截然不同的形式，我仍要毫不松懈地坚持，祂是从无中创造了这些东西，因为在祂产生它们之前，它们根本不存在。那么，如果从未存在的东西成了存在，而从未存在等于无，则从无中产生和从某物中产生有何区别？反之亦然；因为曾经存在等于曾是某物。然而，如果\Emph{有}区别，两种可能性都有利于我的立场。因为如果天主从无中创造了万物，祂就能从无中召出那已化为无的血肉；或者，如果祂从原质中塑造了其他事物，祂也能从别处召出血肉，无论它沉入了何种\Emph{深渊}。无疑，那位创造者最适宜再造，因为产生比复生更伟大，赋予起始比维持延续更伟大。根据这个原则，你可以确信，血肉的复原比它的最初形成更容易。\par

% structure: p0025 source=h2 target=section reason=relative-heading-rank; base=h2

\section{第十二章 自然界中确证血肉复活的一些类比}

试想那些（我们刚才提及的）天主大能的类比。白昼死去变为黑夜，被埋葬于黑暗各处。世界的荣耀昏暗于死亡的阴影中；其整个实质被黑色玷污；万物变得污秽、寂静、愚蠢；各地事务停歇，劳作止息。于是人们为光明的失落而哀悼。然而它又再次复苏，带着自己的美丽、自己的嫁妆、自己的太阳，一如既往，完整无缺，普照世界，杀死自己的死亡——黑夜，打开自己的坟墓——黑暗，作为自己的继承者出来，直到黑夜也复苏——它也带着自己的随从。因为星辰的光芒重新点燃，那些在晨光中熄灭的；星宿的遥远群组再次被带回视野，这些曾被\Emph{白日}的短暂间隔从视线中移除。月亮的镜子也重新装饰，它每月的行程曾将其磨蚀。冬季和夏季归来，春潮和秋日也是如此，带着它们的资源、它们的规律、它们的果实。因为大地从天上接受教导，给剥光的树木穿衣，重新给花朵染色，再次铺展青草，再生产被消耗的种子，而且直到被消耗后才再生产。奇妙的方法！从一个骗取者变成一个保存者，为了恢复而拿走；为了守护而摧毁；为了使之完整而伤害；为了使之扩大而先缩小。（这个过程）确实归还给我们比它剥夺的更丰富更圆满的祝福——通过一种有益的毁灭，一种有利的伤害，一种增益的损失。总而言之，我要说，整个受造界充满更新。你偶然遇到的任何事物，早已存在；你失去的任何事物，必会再次归来。万物在消失后回归原状；万物在结束后开始；它们结束正是为了重新开始。没有事物消灭，只是为了救恩。因此，这整个轮回秩序见证了死者的复活。天主在祂的工程中写下了这一点，早于祂在圣经中写下它；祂在祂大能的作为中宣告了它，早于祂在默感的话语中。祂先差遣自然作为你的老师，意欲再差遣预言作为补充的导师，使你作为自然的学生，更容易相信预言，并在听到你早已在各方面所见的事情时，毫不迟疑地接受（其见证）；也不怀疑你所发现为万物的恢复者的天主，同样是肉身的复活者。而且，既然万物都是为人而复活，为人使用而提供——但不为人，除非也为他的肉身——那么，那肉身本身，为了它并为服事它，没有事物归于虚无，它自己又怎能完全消亡呢？\par

% structure: p0027 source=h2 target=section reason=relative-heading-rank; base=h2

\section{第十三章。从我们的作者对第九十二篇圣咏一节经文的看法，凤凰成为我们身体复活的象征。}

然而，如果整个自然只是模糊地描绘我们的复活；如果受造界没有提供完全相似的记号，因为它的各种现象很难说是死亡，不如说是结束，也不被视为重获生命，而只是重新形成；那么就取一个我们盼望的最完整、最无可辩驳的象征吧，因为它将是一个有生命的存有，同样隶属于生命和死亡。我指的是东方特有的那种鸟，因其独一无二而闻名，因其死后生命而神奇，它以自愿的死亡更新生命；它的死日即是它的生日，因为在那一天它离开又回来；刚才还没有凤凰，如今又是一只；刚才还不存在，如今又重现；另一个，却又是同一个。还有什么比这更明白、更切合我们的主题的呢？或者这样的现象还能见证什么别的事呢？天主甚至在祂自己的圣经中说：\BibleQuote{\Emph{义人}必如凤凰繁茂。}就是说，必从死亡、从坟墓繁茂或复苏——为教导你相信，肉体物质甚至可以从火中恢复。主曾宣告我们\BibleQuote{比许多麻雀更好：}\BibleRef{玛 10:33}好吧，即使不比许多凤凰更好，也算不得什么大事。但人必须一次而永远地死去，而阿拉伯的鸟却确能复活吗？\par

% structure: p0029 source=h2 target=section reason=relative-heading-rank; base=h2

\section{第十四章。一个充足的理由出现在未来对人的审判中，促成了肉身的复活。它将同样关注身体的工作和灵魂的工作。}

既然这乃是天主在\Emph{自然界的}比喻中和在祂亲口所说的话语中所彰显的神圣作为之概要，现在就让我们来接近祂的谕令和法令本身吧，因为这是我们在主题材料中所主要采用的划分。我们开始探讨肉身的尊荣，即它是否具有这样的本性：一旦被毁灭，便有可能恢复原状。然后我们追查关于天主的能力的问题，即这能力是否足够大，足以习惯性地将这种恢复赐予已被毁灭的事物。现在，如果我们已证实了这两点，我希望你接着探询那（\Emph{原因的问题}），即这原因是否足够重大，足以将肉身的复活视为必要，并在各方面都合乎理性；因为在这抗辩背后隐藏着这样的理由：肉身可能完全能够恢复，天主也完全能够实现这恢复，但恢复的原因必须先已存在。那么，你们这些认识一位既是至善又是至义的天主的人，请承认一个充分的原因吧——祂的至善出自祂自己的本性，祂的至义则是因我们而有的。因为，如果人从未犯过罪，他就只会且仅会从其本性之属性上认识天主的至善。但现在，他却因一个原因的需要而经验到祂也是一位至义的天主；不过，就在这同一情形下，祂仍然保留着祂的卓越之善，同时祂也是至义的。因为，祂既扶助善人又惩罚恶人，以此彰显祂的正义，同时使这两种行为都成为祂善的明证：一方面施行惩罚，另一方面赐予赏报。但关于\Person{马尔基昂}，你将有机会更充分地学习这是否就是天主的全部属性。与此同时，我们的天主是如此完美，以致祂是正义的审判者，因为祂是主；祂是主，因为祂是造物主；祂是造物主，因为祂是天主。由此导致，那位我不知其名的异端者认为，祂不是正当的审判者，因为祂不是主；不是正当的主，因为祂不是造物主。因此，我不知道祂如何能是天主，祂既不是造物主（而天主是造物主），又不是主（而造物主是主）。既然如此，对于身为天主、主和造物主的这位\Emph{伟大存有}来说，最合适的莫过于召叫人来接受审判，审判的问题正是：人是否注意承认并尊崇他的主和造物主。这就是复活将要成就的那种审判。因此，复活的全部原因，或者更确切地说，复活的必要性，就在于这一点：那最终审判的安排应最合于天主。现在，在实现这一安排时，你必须考虑天主的审断是否监督着对人的两种本性——灵魂和肉身——的司法审查。因为适合受审的对象，也有资格被复活。我们的立场是，首先必须相信天主的审判是全面的，然后是绝对的，以致是终局的，因此是不可撤销的；也是公义的，不对任何一部分特别严苛；更是相称于天主的，完整而确定，与祂伟大的忍耐保持一致。因此，审判的圆满和完全，仅在于代表整个人的利益。既然整个的人由两种本性的结合构成，他必须在这两种本性中都出现，因为宜于在他整体中受审；当然，他生命的过程中也从未离开过他的整体状态。因此，正如他生活过，他也必须被审判，因为他要就他生活的方式受审。因为生命是审判的原因，它必须在它所曾拥有的、在行使生命机能时的那许多本性中接受调查。\par

% structure: p0031 source=h2 target=section reason=relative-heading-rank; base=h2

\section{第十五章 肉身既在人类一切行为中有分于灵魂，就必在永生的赏报中也有分。}

来吧，让我们的对手在生活的事工上割裂肉身与灵魂的联系，好叫他们敢在生命的报偿上也割裂它。让他们否认它们在行动中的联合，以便他们能公平地否认它们在赏报中的共享。肉身若在原因中无份，就不该在判决中有份。如果灵魂独自离去，就让它独自被召回。但（从未发生过这样的事）；因为灵魂独自并未离开生命，就像它独自跑完它所离开的历程——我是指今生。的确，灵魂独自远离经营生活，以至于我们连思念也不从与肉身的共融中撤出，无论这些思念多么孤立，多么未通过肉身而化为行动；因为凡在人心内所做的，都是由灵魂在肉身内、偕同肉身、并通过肉身而做的。总之，主责备我们的思念时，也将肉身的这一方面（人心），灵魂的堡垒，包括在祂的谴责中：\BibleQuote{你们为什么心里思念恶事？}\BibleRef{玛 9:4}又说：\BibleQuote{凡注视妇女，有意贪恋她的，他已在心里奸淫了她。}\BibleRef{玛 5:28}因此，即使没有运作、没有效果的思念，也是肉身的行为。但若你承认那主宰感官的官能——他们称之为\OriginalTerm{主导官能}{Hegemonikon}——其圣所位于大脑、或眉间、或哲学家乐于放置的任何地方，肉身仍是灵魂的思维场所。灵魂在肉身内时，从未离肉身。肉身不与灵魂共同办理的事，一件也没有，因为离开灵魂它就不存在。也请仔细考虑，思念是否不由肉身管理，因为正是透过肉身，思念才在外部被辨别和认识。让灵魂刚筹划一个计谋，面容就露出来——面容是我们一切意向的镜子。他们可以否认行动中的结合，但不能否认念头中的协作。然而他们列举肉身的\Emph{罪过}；那么为它的犯罪行径，它必被归于惩罚。但我们还要对他们列举肉身的\Emph{美德}；那么为它的德行行为，它也当得未来的赏报。再者，既然在一切事上都是灵魂行动并推动我们，那么肉身的职能就是服从。现在我们不得认为天主是不义或怠惰的。不义（无论如何，祂是不义的），若祂将善工中有份的\Emph{肉身}排除于赏报之外；怠惰，若祂免除它在恶行中的共犯的惩罚：而人的审判被认为更完善，当它在每一行为中发现代理人，既不饶恕\Emph{有罪的}，也不吝惜\Emph{有德的}，使他们在惩罚或赞美中与雇用他们服务的主子分享全份。\par

% structure: p0033 source=h2 target=section reason=relative-heading-rank; base=h2

\section{第十六章 异端者称肉身是\InnerQuote{灵魂的器皿}，以摧毁肉身的责任。他们的诡辩反噬自身，显示肉身是人的行动的同工。}

然而，当我们把权威归于灵魂，把服从归于肉身时，我们必须注意，不要让（我们的对手）用另一个论点来推翻我们的立场，他们坚持将肉身置于对灵魂的服务中，以致它不被（视为）灵魂的仆人，免得他们在这样看待的情况下，被迫承认它（与灵魂）的同伴关系。因为他们会争辩说，仆人和同伴在履行各自职责时拥有判断力，并且在两种关系中都有意志的能力；简而言之，（他们会声称自己是）人本身，因此（会期望）与他们的主人分享功劳，他们自愿将帮助交给主人；然而，肉身本身没有判断力，没有感情，而且没有自己意志或拒绝的能力，事实上，它对灵魂而言似乎只是一个器皿，而不是一个仆人。只有灵魂因此将要在（最后审判时）首要地被审判，关于它如何使用肉身这器皿；器皿本身当然不能接受司法裁决：因为如果有人把毒药混进杯里，谁会定这杯的罪？或者如果有人用剑犯了强盗的暴行，谁会判剑被野兽吞噬？好了，现在，我们承认肉身是无辜的，只要恶行不归于它：那么，有什么能妨碍它因无罪而得救呢？因为虽然它没有善行的功绩，也没有恶行的罪过，但拯救无辜者更符合天主的良善。确实，一个慈善的人有义务这样做：那么，适合至慷慨者性情的，就是甚至无偿地赐予这样的恩惠。然而，至于杯子，我不取那被注入某种致命毒药的杯子，而是取那被淫荡女人的气息，或库柏勒司祭的，或角斗士的，或刽子手的气息所污染的杯子：那么，我想知道，你是否会对这个杯子判处比这些人亲吻更轻的刑罚？一个被我们自己的污秽玷污的杯子，或者一个不是按我们心意调和的杯子，我们往往会把它摔碎，然后把怒气发泄在我们的仆人身上。至于那把沾满强盗受害者鲜血的剑，谁不会把它完全赶出家门，更不用说卧室或枕头边，因为可以推定他一定会梦见那些灵魂的显现，这些灵魂因为他枕着那曾经流出它们自己鲜血的刀刃而追逐并困扰着他？然而，拿那只没有瑕疵的杯子，它值得忠信服务的称赞，它的主人会用花冠装饰它，或用一把鲜花来荣耀它。那把在战争中染上光荣污迹、参与更好杀人的剑，也会通过奉献获得自己的赞美。那么，甚至对器皿和工具也可以作出决定性的判决，以便它们也能分享其所有者和使用者的功绩。\Emph{我这样说这么多}，是因为渴望回应这个论点，尽管这个例子有不足之处，因为对象的性质不同。因为每个器皿或每个工具都是从外部变得有用，其材料完全外在于其人类\Emph{主人或使用者}的本质；而肉身，从它在母腹中的最初存在起就与灵魂一同被孕育、形成和诞生，也在它的一切行动中与灵魂混合。因为虽然它被宗徒称为\Quote{器皿}，正如他命令人要\Quote{以尊重}对待它，\BibleRef{得前 4:4}，但同一位宗徒也称之为\Quote{外在的人}，\BibleRef{格后 4:16}——当然，那泥土起初就被赋予了人的称号，而不是杯子或剑或任何无价值的器皿。现在，它之所以被称为\Quote{\Emph{器皿}}，是考虑到它的容量，藉以接受和容纳灵魂；但称之为\Quote{\Emph{人}}，则是基于本质相通，这使它在一切行动中成为仆人而非工具。因此，在审判中，它将被视为仆人（即使它本身没有独立的判断力），理由是它是拥有这种判断力的那个实体的组成部分，而不只是动产。虽然宗徒很清楚，肉身自己什么也不做，除非这也要归于灵魂，但他仍然认为肉身是\Quote{\Emph{有罪的}}；\BibleRef{罗 8:3}，免得它因为似乎被灵魂推动而误以为完全不负责任。同样，当他把某些值得称赞的行为归于肉身时，他说：\Quote{\Emph{所以}要在你们身体上光荣和显扬天主}，\BibleRef{格前 6:20}——他确信这些努力是由灵魂驱动的；但他仍然把它们归于肉身，因为他也应许赏报给肉身。此外，如果肉身是无可指摘的，责备就不适合它；如果它不能承受光荣，劝勉也不适合它。事实上，责备和劝勉对肉身都是空的，如果它是那种在复活中确实会得到赏报的不当对象的话。\par

% structure: p0035 source=h2 target=section reason=relative-heading-rank; base=h2

\section{第十七章 肉身将与灵魂一同承受最终审判的惩罚判决。}

每一个同意我们意见的无学问之人，可能会倾向于认为，甚至仅基于这个理由，肉身必须出现在\Emph{最终}的审判中，因为否则灵魂作为无形体的东西，将无法承受痛苦或快乐——这是普遍的见解。然而，我们这边却在此主张，并在关于这个主题的一篇专门论文中证明，灵魂是有形体的，在它的本性中拥有一种特殊的坚固性，使之既能感知也能承受。即使灵魂在阴间脱离了肉身，且与肉身分离，却仍然能受痛苦和享祝福，这由\Person{拉匝禄}的例子所证明。我无疑给我的对手留下了话柄说：既然灵魂自身有一个有形体的实体，它便足够拥有承受和感觉的能力，而无需肉身的在场。不，不（这是我的回答）：它仍然需要肉身；不是因为它没有肉身帮助就不能感觉任何东西，而是因为它必须与肉身一起拥有这种能力。因为\Emph{就}它自有充足的行动力\Emph{而言}，它\Emph{因此}也拥有承受的能力。但事实是，就行动而言，它多少有些无能；因为就其本性而言，它仅仅具有思考、意愿、欲望、处置的能力；为了完全实现目的，它寻求肉身的协助。同样，它也需要与肉身的结合来承受苦难，以便通过肉身的帮助，它能像没有协助时无法完全行动那样，完全地承受苦难。确实，对于那些它自己就有能力犯的罪，如贪欲、思想和意愿，它立即为它们付出代价。现在，毫无疑问，如果这些本身足以构成绝对的功过，而不需要加上\Emph{行为}，那么灵魂独自就足以面对审判的全部责任，因为它独自就有足够的能力去做这些事情而被审判。然而，既然行为也不可分割地附属于功过；既然行为也是由肉身执行完成的，那么灵魂脱离肉身后，就不足以因那些实际上是肉身的工作而蒙受快乐或痛苦的报应，即使它有自己的身体，有自己的肢体，这些同样不足以使它完全感知，正如它们也不足以使它完全行动一样。因此，就像它在每个个别情况中行动一样，它在阴间也相应地承受，最先品尝审判，正如它最先怂恿人犯罪；但它仍在等待肉身，以便通过肉身也为它的行为做出补偿，因为它把自身思想的执行加给了肉身。简而言之，这将是被推迟到最后大日子的审判过程，以便通过肉身的显现，天主的全部报复得以完成。此外，（显然可以看出）如果那已经由灵魂在阴间品尝的审判注定只针对灵魂，就不会拖延到结局。\par

% structure: p0037 source=h2 target=section reason=relative-heading-rank; base=h2

\section{第十八章 圣经的语句和段落明确宣告\Quote{死者的复活}。这个短语本身的力度解释了肉身在普遍复活中的突出位置。}

迄今为止，我通过前言旨在为所有许诺肉身复活的圣经辩护奠定基础。现在，既然\Emph{这一真理}得到了许多公正合理的考量支持——我指的是肉身自身的尊严，天主的权能与大力，彰显这些的类比情形，以及审判的充分理由及其必要性——那么，圣经理所当然应当按照这些权威性考量所提示的意义来理解，而不是按照异端者的妄念，这些妄念完全源于不信，因为肉身从死亡中复原并恢复生命被视为不可信；并非因为（这种复原）是肉身本身无法达到的，或是天主无法做到的，或是不适合\Emph{最终的}审判。无疑，如果它未曾在天主的话中启示，它可能是不可信的；不过，即使天主未曾如此预先宣布，我们也可以相当合理地假定，它的启示之所以被保留，仅仅是因为已经提供了许多有利于它的强有力的推定。然而，既然（这一伟大事实）在如此多受默感的段落中被宣告，这本身就有力地劝阻我们不要以一种与这些劝服我们接受它的论证所证实的意义不同的方式来理解它，即使不考虑启示的见证。那么，首先让我们看看，我们的这一希望是以什么标题呈现在我们眼前的。我想，有一道神圣谕令展示在众人眼前：\Quote{死者的复活}。这些话语简洁、果断、清晰。我打算仔细检查这些词语，讨论它们，并发现它们适用于什么实体。至于\OriginalTerm{复活}{resurrectio}一词，每当我听说它即将临到一个人身上时，我不得不追问他的哪一部分被注定要\Emph{倒下}，因为除非先已倒下，否则无可期望重新起来。只有那些不知肉身因死亡而倒下的人，才无法发现肉身因生命而站立。大自然宣告了天主的判决：\BibleQuote{你原是尘土，还要归于尘土。}\BibleRef{创 3:19} 即使没有听过这一判决的人，也看到了事实。没有死亡不是我们肢体的毁灭。主也描述了身体的这一命运，当祂披戴着同一实体（即肉身）时，祂说：\BibleQuote{你们拆毁这座圣殿，三天之内我要把它重建起来。}\BibleRef{若 2:19} 因为祂表明，被拆毁、被推倒、被压制属于什么——即属于那也要被举起和复起的东西；虽然祂同时带着一个\BibleQuote{忧闷至死的灵魂}\BibleRef{玛 26:38}，但那灵魂并没有因死亡而倒下，因为圣经告诉我们：\BibleQuote{祂是论及祂的身体。}\BibleRef{若 2:21} 所以，是肉身因死亡而倒下；因此，它从\OriginalTerm{倒下}{cadendo}而得名\OriginalTerm{尸体}{cadaver}。然而，灵魂在其名称中毫无\Emph{倒下}的痕迹，因为它的状态中确实没有可死性。不，正是灵魂在离开身体时将其毁灭传递给身体，正如它将在重新进入身体时再次将其从地上举起。凡是因进入而举起的不可能倒下；凡是因离去而造成毁灭的不可能衰颓。我更进一步说，灵魂甚至不与身体一同入睡，也不与它的同伴安卧休息。因为它在梦中激动不安；然而，如果它躺下，它本可安息；如果它倒下了，它当然会躺下。因此，那甚至不陷入死亡相似之物的，也不会屈服于死亡的实况。现在转到\Emph{另一个}词\OriginalTerm{死者}{mortuorum}，我希望你仔细察看，看看它适用于什么实体。如果我们在这点上同意异端者有时所持的观点，即灵魂是可死的，那么既然可死，它将要获得复活；这就会提供一种推定，即同样可死的肉身也将分享\Emph{同一}复活。但我们目前的重点是，从这个词的适当含义中得出它所指示的命运观念。现在，正如\Emph{复活}一词用于那倒下之物——即肉身——同样，\Emph{死的}一词也有相同的应用，因为所谓\Quote{死者的复活}指的是那已倒下之物的再次起来。我们从信德之父\Person{亚巴郎}的事例中学到这一点，他是与天主亲密交往的人。因为他向赫特人请求一块地方埋葬\Person{撒辣}时，对他们说：\BibleQuote{求你们给我一块坟地，好让我埋葬我的死者，}\BibleRef{创 23:4}——这当然是指她的肉身；因为他不可能渴望一个地方来埋葬她的灵魂，即使灵魂被视为可死的，即使它能被\Quote{ \Emph{死的}}这个词来描述。所以，既然这个词指示的是身体，那么当谈到\Quote{死者的复活}时，所指的就是\Emph{人}的身体的再次起来。\par

% structure: p0039 source=h2 target=section reason=relative-heading-rank; base=h2

\section{第十九章　异端者对\Quote{死者复活}这一短语所赋予的诡辩式理解，仿佛它意指新生命的道德转变。}

现在，对所述词组及其含义的考量——当然，除了维护重要词汇的真实意义之外——必然有助于进一步的结果：无论我们的对手以比喻和寓意语言为借口，在主题上投下何等模糊，真理将在更清晰的光中显现，并从不确定性中规定出确定而明确的规则。因为有些人，当他们碰到一种非常常见的先知性陈述形式，通常以比喻和寓意表达，尽管并非总是如此，甚至将论述最清楚的死者复活道理曲解为某种想象的意义，声称连死亡本身也必须按精神意义理解。他们说，通常被认为是死亡的——即灵魂与身体的分离——并非真正如此；倒不如说是不认识天主，人因此而死于天主，且不亚于被埋葬在坟墓中而埋葬在错误中。因此，那也必须被当作复活，即一个人通过接近真理而重新获得生命，驱散了无知之死，并藉天主获得新生命，从旧人的坟墓中迸发出来，正如主曾将经师和法利塞人比作\Quote{粉饰的坟墓}。\BibleRef{玛 23:27}由此推论，那些藉信德达到复活的人，一旦在圣洗中穿上主，便与主同在。因此，他们常以这种狡猾，甚至在谈话中误导我们的弟兄，仿佛他们也\Emph{和我们一样}持守死者复活。他们说：\Quote{祸哉，那在现世身体中没有复活的人！}因为他们担心，如果立即否认复活，会吓倒听众。然而，在他们心里却秘密这样想：祸哉，那个在今生未能发现异端奥秘的傻瓜！因为在他们看来，这就是复活。然而，也有许多人，声称持守灵魂离世后的复活，主张走出坟墓就是脱离世界，因为他们认为世界是死人的居所——即那些不认识天主的人；他们甚至会走得更远，说这实际上意味着脱离身体本身，因为他们想象身体拘禁灵魂，当灵魂被关闭在世俗生命的死亡中时，如同在坟墓中。\par

% structure: p0041 source=h2 target=section reason=relative-heading-rank; base=h2

\section{第二十章 比喻的意义基于字面事实。此外，比喻风格并非先知经书中唯一的一种，如异端者所声称的。}

现在，为推翻所有这些自负之见，让我立刻驱散它们所依据的初步观念——他们声称先知们所有的宣告都是以比喻方式作出的。倘若真是如此，比喻本身就不可能被区分，因为比喻所基于的真理并未被宣告。而且，如果一切都是比喻，那比喻所指的对象又在何处？如果没有脸的存在，你如何能举起镜子照你的脸？但事实上，并非一切都是比喻，也有字面陈述；也不是一切影子，也有实体本身。因此，我们关于主本身的预言甚至比白昼还要清晰。因为童贞女在子宫中怀孕并非比喻；她生厄玛努尔，即耶稣，天主与我们同在，也非借喻。即使祂是比喻性地夺取大马士革的能力和撒玛利亚的战利品，\BibleRef{依 8:4}但祂要\Quote{与民间的长老和首领进入审判}却是字面意义。\BibleRef{依 3:13}因为在比拉多身上，\Quote{外邦人喧嚷}；在以色列身上，\Quote{人民谋算虚妄的事}；\Quote{地上的君王}在黑落德身上，\Quote{首领}在亚纳斯和盖法身上，一同聚集\Quote{反对主，反对祂的受傅者}。祂又被\Quote{牵到宰杀之地，又像羊在剪毛的人面前}，即黑落德，\Quote{哑口无声，祂竟不开口}。\BibleRef{依 53:7}\Quote{祂将背转给鞭打，将面颊转给掌击，连被唾弃的羞辱也不转脸}。\Quote{祂被列在罪犯之中}；\BibleRef{依 53:12}\Quote{祂的手和脚被刺透}；\Quote{他们为祂的衣服拈阄}；\Quote{他们给祂苦胆，又拿醋给祂喝}；\Quote{他们摇头讥笑祂}；\Quote{祂被叛徒估价三十块银钱}。\BibleRef{匝 11:12}什么比喻？依撒意亚在这里给了我们什么比喻？达味给了什么修辞？耶肋米亚给了什么寓言？甚至祂的奇事，他们也没有用比喻的语言描述。不然，盲人的眼睛岂没有睁开？哑巴的舌头岂不恢复说话？\BibleRef{依 35:5}松懈的手和瘫痪的膝岂不变得强壮，跛子岂不像鹿一样跳跃？无疑，我们也习惯按主所医治的身体疾病的类比，给这些预言陈述赋予精神意义；但它们都已字面应验，从而表明先知预见了两种意义，只是他们许多的话只能以纯朴简单的含义理解，毫无比喻的晦涩，如我们听到关于列国和城市的倾覆，关于提洛、埃及、巴比伦、厄东，以及迦太基的舰队；又如他们预言以色列自身的惩罚和宽恕、它的充军、复兴，以及最后的最终分散。谁会宁愿给所有这些事件附上隐喻的解释，而不接受它们的字面真实？现实蕴含在话语中，正如话语在现实中读到。因此，我们发现比喻风格并非出现在先知记录的所有部分，尽管它偶尔出现在某些段落中。\par

% structure: p0043 source=h2 target=section reason=relative-heading-rank; base=h2

\section{第二十一章 在\Quote{死者复活}这一词组中绝非隐喻。永恒真理的重要性越高，其圣经宣告的清晰度也越高。}

好吧，如果你说，在圣经的某些部分偶尔会出现这种情况，那么，为什么在那句话中，复活可能被灵意理解时，却不是这样呢？有多个理由说明为什么不是。首先，圣经中那么多重要段落，它们如此明显地证实了身体的复活，以至于连比喻性意义的表象都不容许，那它们的意义必须是什么？确实，（因为有些段落比其他段落更为隐晦），正如我们上面所指出的，正确的做法必然是：不确定的陈述应由确定的陈述来决定，隐晦的应由清晰明白的来决定；否则，恐怕在确定与不确定、明确与隐晦的冲突中，信德会被粉碎，真理会陷入危险，而天主本身也会被指责为反复无常。于是便产生了这样的不可能性：我们的整个信赖完全依赖于其上，我们的教导也完全基于其上的这个奥秘本身，竟然会显得是被模棱两可地宣布和隐晦地提出的，因为复活的希望，除非在惩罚和奖赏两方面都清楚地阐明，否则就无法说服任何人接受像我们这样的宗教——它暴露于公众的憎恶和被指责为敌视他人的指控之下。没有确定的工作，报酬就不确定。没有真正的恐惧，危险就只是可疑。但是，无论是奖赏的回报，还是失去它的危险，都取决于复活的结果。现在，即使天主对那些城市、国家和君王的旨意——它们仅仅是暂时的、地方性的和个人的——在预言中都已经如此清楚地宣告，那么，祂那些永恒的和关乎全人类普遍利益的安排，又怎能被认为在其自身中缺乏一切真正的光明呢？这些安排越伟大，它们的宣布就应该越清晰，以便它们更伟大的卓越性能被相信。我认为，天主绝不可能被归因于嫉妒、诡诈、不一致或狡诈——正是借助这些邪恶的品质，所有非凡宏大的计划才得以争议性地颁布。\par

% structure: p0045 source=h2 target=section reason=relative-heading-rank; base=h2

\section{第二十二章 圣经禁止我们假设复活已经过去，或在死亡时立即发生。我们的希望和祈祷指向最后的大日作为它完成的时期。}

在这一切之后，我们还必须留意那些禁止我们相信像你们的\OriginalTerm{动物派者}{Animalists}（因为我不愿称他们为\OriginalTerm{精神派者}{Spiritualists}）所持守的复活的圣经章节；这种复活要么被认为在人一获得真理的知识时就\Emph{发生}，要么在离世后即刻完成。既然我们全部希望的时节已在圣经中确定了，并且按照我的理解，我们不允许将其成就置于基督来临之前，我们的祈祷便指向今世的终结，即它在主的大日——祂愤怒与报复的日子——的消逝；那最后的日子，隐藏（于众人），除了圣父无人知晓，尽管已借征兆、奇事、元素的解体及万国的冲突预先宣告了。如果主自己什么也没说（尽管预言本就是主的话），我就会引用先知的话；但祂亲口证实了\Emph{他们的陈述}，这更符合我的目的。当祂的门徒问祂，祂刚刚所说的关于圣殿毁灭的那些事何时发生时，祂先向他们讲述了犹太事件的次序，直到耶路撒冷倾覆，然后讲述了关乎万民直到世界末了的事。因为祂在宣告了\Quote{耶路撒冷要被外邦人践踏，直到外邦人的时期满了，}\BibleRef{路 21:24}——当然指的是那些被天主拣选并与以色列遗民一同聚集的人——之后，接着向这世界和时代宣告（正如岳厄尔、达内尔和众先知一致同意的）：\BibleQuote{在太阳、月亮和星辰上将有征兆，万民因困惑而困苦，海洋和波涛咆哮，人因恐惧和期待即将临到地上的事而心灰意冷}\BibleRef{路 21:25}。祂说：\BibleQuote{因为天上的权势将要动摇；那时他们要看见人子带着大能力和大荣耀驾云降临。当这些事开始发生时，你们要挺身昂首，因为你们的救赎近了。}祂说的是救赎\Quote{近了}，而不是已经临在；又说是\BibleQuote{那些事开始发生}，而不是已经发生：因为当那些事发生时，我们的救赎就临近了，直到那时，救赎一直被说成是临近的，以提升并激动我们的心，指向那时我们希望的临近收获。祂随即附上了一个关于此事的比喻：\BibleQuote{树木嫩芽萌发成花梗，然后开花，那是果实的前兆。}祂接着说：\BibleQuote{同样，当你们看见这一切事发生时，就知道天国近了。}\BibleQuote{所以，你们要时时警醒，时常\Emph{祈祷}，使你们能配得逃脱这一切事，并站立在人子面前；}\BibleRef{路 21:36}这无疑是指在复活时，在这一切事先发生之后。因此，虽然在认识这全部奥秘方面已有萌芽，但只有在主亲自临在时，花才绽放、果才结出。那么，是谁如此不合时宜、如此严厉地惊动了现在在天主右手的\Quote{主}，以至于祂\Quote{剧烈震动}（如依撒意亚所言，\BibleRef{依 2:19}）\Quote{大地}？——我以为是至今未碎的大地。谁如此早地将\BibleQuote{基督的仇敌放在祂脚下}（借用达味的话），使祂比圣父更急切，而我们集会中的每一群人仍在喊着\Quote{把基督徒丢给狮子}？有谁曾看见耶稣从天降下，如同宗徒们看见祂升天一样，正如那\Emph{两位}天使所指定的？\BibleRef{宗 1:11}直到如今，他们尚未按支派捶胸痛哭，仰望他们所刺透的那位。还没有人遇见厄里亚；\BibleRef{拉 4:5}没有人逃脱假基督；\BibleRef{若一 4:3}没有人必须哀哭巴比伦的倾覆。\BibleRef{默 18:2}那么现在有谁复活了呢？除非是异端者。当然，\Emph{他}早已离开自己尸体的坟墓——尽管他至今仍易患热病和溃疡；他也已践踏了仇敌——尽管他至今仍要与世界权势争斗。而他当然已是君王——尽管他至今仍将凯撒的归给凯撒。\BibleRef{玛 22:21}\par

% structure: p0047 source=h2 target=section reason=relative-heading-rank; base=h2

\section{第23章 圣保禄关于属灵复活的若干段落，与身体将来的复活相符，且已为经文所预设}

宗徒在\WorkTitle{致哥罗森人书}中确实教导说，我们从前是死的，在我们的心思中与主疏远、为敌，在我们生活在恶行中的时候；\BibleRef{哥 1:21}并且我们曾藉着圣洗与基督一同埋葬，也藉着使祂从死者中复活的天主的运行所赐的信德，与祂一同复活。\BibleRef{哥 2:12}\Quote{你们（他接着说）从前因罪恶和未受割损的肉体是死的，天主却使你们与祂一同生活，赦免了你们的一切过犯。}又说：\Quote{如果你们同基督一起脱离了世俗的要素，为什么还像在世界上生活一样，服从规条呢？}既然他使我们在灵性上死亡——但却是以这样一种方式，即允许我们有一天必须经历身体死亡——同样，他既认为我们也已在类似灵性意义上复活，也同等允许我们将进一步经历身体的复活。他明说：\Quote{既然你们与基督一同复活了，就该追求天上的事，那里有基督坐在天主的右边。你们要思念天上的事，不要思念地上的事。}\BibleRef{哥 3:1}因此，他指出我们（与基督）复活是在我们的心思中，因为唯有藉此我们目前才能企及天上的事物。这些事物我们不应\Quote{追求}，也不应\Quote{思念}，如果我们已经拥有它们的话。他接着又说：\Quote{因为你们已经死了}——他意指死于罪恶，而非死于自己——\Quote{你们的生命已与基督一同藏在天主内。}那隐藏的生命尚未被掌握。同样，\Person{若望}说：\Quote{现在还不显明我们将成为什么，但我们知道，当祂显现时，我们将相似祂。}\BibleRef{若一 3:2}我们远未成为我们所不知道的；如果我们已经是（相似祂），我们当然就会知道。因此，他所谈的是在今世因信德而默观我们蒙福的希望——不是其临在或拥有，而只是其期待。关于这期待与希望，\Person{保禄}写信给迦拉达人：\Quote{我们却是因圣神，由信德期待正义的希望。}\BibleRef{迦 5:5}他说\Quote{我们期待它}，而非我们拥有它。他指的天主的义，就是我们将因我们\Emph{行为}的赏报而受的审判。正是为了这期待，宗徒写信给斐理伯人：\Quote{\Emph{或许}我能达到从死者中的复活。这不是说我已经达到，或已经成全了。}\BibleRef{斐 3:11}然而，他已经相信并认识了一切奥秘，是特选的器皿和外邦人的\Emph{伟大}导师；尽管如此，他接着说：\Quote{我只顾向前追赶，或许我能获得，因为我也已被基督耶稣获得。}不仅如此：\Quote{弟兄们，（他接着说道）我不以为我已经获得了；我只顾一件事：忘记背后的，努力面前的，向着标竿直跑，为要获得无可指摘的奖赏，藉此我能达到它；}意指在适当的时候从死者中的复活。正如他对迦拉达人所言：\Quote{我们行善不要厌倦：因为\Emph{到了一定的时期}，我们必要收获。}\BibleRef{迦 6:9}同样，关于\Person{敖乃息佛洛}，他也写信给\Person{弟茂德}：\Quote{愿主在那一天赐他获得仁慈；}\BibleRef{弟后 1:18}正是针对那一天和那个时辰，他嘱咐弟茂德本人\Quote{要持守所托付给你的，毫无玷污，无可指摘，直到我们的主耶稣基督的显现：到了日期，那蒙福的、独一的掌权者，万王之王、万主之主要显示它}——论及（祂是）天主。正是这些时期，\Person{伯多禄}在\WorkTitle{宗徒大事录}中提及，他说：\Quote{所以你们应当悔改，并归正，好使你们的罪过得以消除，并使安乐的时期从主面前来到；祂要派遣耶稣基督，就是预先给你们宣讲的：基督必须留在天上，直到万物复兴的时期，这时期天主藉着祂的圣先知的口早已说过。}\BibleRef{宗 3:19}\par

% structure: p0049 source=h2 target=section reason=relative-heading-rank; base=h2

\section{第二十四章 引自圣保禄的其他经节，明确断言肉体在最后审判时的复活}

你们要与得撒洛尼人一同学习这些时期的特征。因为我们读道：\BibleQuote{你们怎样离开了偶像，皈依了天主，来服事永生的真天主，并期待他的圣子从天降下，就是天主从死者中所复活的那一位，耶稣。} \BibleRef{得前 1:9} 又说：\BibleQuote{我们的希望、喜乐、或称耀的冠冕是什么呢？不就是在我们的主天主耶稣基督来临时，你们也在他面前吗？} 同样：\BibleQuote{在我们的天主父面前，在我们主耶稣基督来临时，与他的全体圣人一起。} \BibleRef{得前 3:13} 他教导他们必须\Quote{不要为那些睡了的人忧伤}，同时向他们解释复活的时间，说：\BibleQuote{因为我们如果相信耶稣死了，也复活了，那么藉着耶稣，那些睡了的人，天主也必与他们一同带来。我们凭主的话对你们说：我们这些活着存留到主来临的人，决不会在那些已睡了的人以前。因为主必亲自从天降来，有呼声、有总领天使的声音、有天主的号角；那在基督内死了的人要先复活。然后我们这些活着还存留的人，要同时与他们一起被提到云中，去迎接主在空中。这样我们就常与主在一起。} \BibleRef{得前 4:13} 如今有何总领天使的声音（我想问）、有何天主的号角可以听到，难道不是在异端者的娱乐中吗？因为，即便福音的话可以被称为\Quote{天主的号角}，因为它仍在召唤人，然而他们那时必须要么身体死了，才能复活；但这样他们如何活着？要么被提到云中，但这样他们如何在这里？无疑，正如宗徒所说，他们\Quote{最可怜的}是\Quote{只在今生}有盼望的人：\BibleRef{格前 15:19} 他们因过早抓住后世应许之事而被排除；在真理上犯错，不亚于\Person{斐理托}和\Person{赫摩革乃}。\BibleRef{弟后 1:15} 因此，圣神以其伟大，清晰地预见到所有这些解释，就在他（宗徒）致得撒洛尼人的这封书信中，提示\Emph{如下}：\BibleQuote{弟兄们，至于时间和季节，没有必要给你们写什么。你们自己清楚地知道，主的日子要像夜间的贼一样来到。当他们说\InnerQuote{平安}和\InnerQuote{一切都安全}时，毁灭就会突然临到他们。} \BibleRef{得前 5:1} 再者，在第二封\Emph{书信}中，他以更大的热忱对他们说：\BibleQuote{弟兄们，关于我们主耶稣基督的来临，和我们对他的聚集，我们恳求你们，不要轻易动摇心智，也不要因神恩、或因言语、或因似乎由我们而来的书信，感到困扰，好像主的日子已经迫近。} 就是\Emph{言语}是假先知的话，\BibleQuote{或因书信}，就是\Emph{书信}是假宗徒的，\BibleQuote{好像从我们而来，以为主的日子已经迫近。不要让人以任何方式欺骗你们。\Emph{因为那日子不会来到}，除非先有叛离之事，} 他指的是这个帝国（罗马）的叛离，\BibleQuote{并且那\OriginalTerm{罪恶之人}{man of sin}被显露出来，} 就是说，\OriginalTerm{假基督}{Antichrist}，\BibleQuote{那\OriginalTerm{丧亡之子}{son of perdition}，他反对并高举自己超过一切称为神或宗教的，以致他坐在天主的圣殿里，宣称自己是神。你们不记得，当我还在你们那里时，我曾对你们说过这些事吗？现在你们知道什么在阻止他，使他在自己的时候被显露。因为\OriginalTerm{邪恶的奥秘}{mystery of iniquity}已经运行，只是那阻止者要阻止，直到他离开。} \BibleRef{得后 2:1} 除了罗马国之外，还有什么障碍呢？它的分裂，因被分散成十个王国，将在（它自己的废墟上）导致假基督出现？\BibleQuote{那时那不法者将被显露，主必以口中的神气将他消灭，并以显现的光荣摧毁他：\Emph{就是那一位}，他的来临是出于撒但对所有能力、奇迹、和谎言的奇妙，以及行将丧亡者在一切不义的欺骗中。} \BibleRef{得后 2:8}\par

% structure: p0051 source=h2 target=section reason=relative-heading-rank; base=h2

\section{第25章 \Person{圣若望}在\WorkTitle{默示录}中同样明确地主张同一伟大教义。}

再者，在\BibleBook{若望}{默示录}中，这些时期的次序展现出来，\Quote{殉道者的灵魂}——他们急切祈求伸冤和审判——在祭台下被教导等候这个次序（\BibleRef{默 6:9}）（我再说，被教导等候），为叫世界首先从天使的盂爵中喝尽等待它的灾祸，\EditorialNote{默示录第十六章}并叫那淫乱之城从十位君主手中接受它应得的惩罚，\EditorialNote{默示录第十八章}叫假基督的兽同他的假先知对天主的教会发动战争；并且，在魔鬼被暂时投入深渊之后（\BibleRef{默 20:2}），初次复活的殊荣将从宝座中颁赐；然后，在他被投入火中之后，最终且普遍复活的审判将依据书卷决定。既然圣经既指出末期的各阶段，又将基督徒希望的圆满实现集中于世界终末本身，那么显然，要么天主所应许我们的一切都在那时得到完成，如此异端者在此处妄称的复活便不攻自破；要么，即使承认对（神圣真理的）奥秘的宣认是一种复活，但在此观点不受损害的情况下，仍有空间相信那为终末所宣告的复活。此外，维持这种属灵的复活本身，就构成了有利于另一种身体的复活的推定；因为倘若那时没有宣告任何复活，就有正当理由只主张这种纯属灵的复活。然而，既然（一种复活）是在最后时刻被宣告的，它就被证明是身体的复活，因为那时也没有宣告属灵的复活。为何只宣告一种性质的复活——即属灵的复活——两次呢？既然这一复活要么应在现在完成，而不顾不同时期，要么就在那时，在所有时期的终末完成。因此，我们这些承认属灵复活在世界终末完全实现的人，更有理由在\Emph{信仰生活}的开始就坚持属灵的复活。\par

% structure: p0053 source=h2 target=section reason=relative-heading-rank; base=h2

\section{第二十六章 就连圣经中关于此主题的比喻性描述也指向身体的复活，这是确保其连贯性与尊严的唯一意义。}

对于前面一个反对意见，即圣经是寓言的，我仍有一个回答要做——我们也可以藉着先知的语言来辩护复活的肉身特性，这些语言同样是比喻性的。请思考天主在称呼人为\Emph{\OriginalTerm{地}{terra}}时所说的那句原始判词：\BibleQuote{你是地，仍要归于地。}\BibleRef{创 3:19}当然，就其肉身的实质而言——这实质是从土中取出的，并且最先接受了人的称号，正如我们已指出的——这段经文岂不也教导我们，要将天主对地所决定的任何义怒或恩宠同样解释为与\Emph{\OriginalTerm{肉身}{caro}}有关，因为严格来说，地并不受祂审判，因为它从未行过任何善恶？\Quote{诅咒，}无疑，它受了诅咒，因为它饮了\Emph{人}的血；\BibleRef{创 4:11}但即便这也是杀人之肉身的预象。因为如果地必须承受喜乐或伤害，那仅仅是为了人的缘故，好使\Emph{他}藉着其居所发生的事件来感受喜乐或悲伤，因此他反而必须付出惩罚，而这惩罚仅仅为了他，连地也必须承受。因此，当天主甚至威胁地时，我宁愿说是祂威胁肉身；同样，当祂向地作出许诺时，我宁愿理解为祂向肉身许诺；正如\Person{达味}的那段经文：\Quote{上主为王，愿地喜乐}——意指圣人们的肉身，这肉身享有天主之国的福乐。随后他又说：\Quote{地看见就战栗；山岭在主的面前如同蜡融化}——无疑是指恶人的肉身；并且（在类似意义上）经上记载：\BibleQuote{他们要仰望他们所刺透的那一位。}\BibleRef{匝 12:10}若有人真以为这两段经文仅仅是指地这一元素而说的，那么它先在主的王权面前欢欣，随后又在主的面前战栗融化，这怎能一致呢？同样在\Person{依撒意亚}中：\BibleQuote{你们必吃地上的美物}\BibleRef{依 1:19}，这话是指当肉身在天主之国中得以更新、相似天使、并等待获得那\BibleQuote{眼所未见，耳所未闻，人心所未想到的}\BibleRef{格前 2:9}事物时所等待的福乐。否则，天主藉田间的果实和今生的元素来邀请人顺从，这何其虚妄，因为祂甚至将这些赐予不虔敬的人和亵渎者；这是基于一次而永远向人颁布的一般条件：\BibleQuote{降雨给义人，也给不义的人；使太阳照耀义人，也照耀不义的人！}\BibleRef{玛 5:45}无疑，信德是有福的，如果它要获取那些天主和基督的敌人不仅使用而且甚至滥用的恩赐，\BibleQuote{崇拜受造物而非造物主！}\BibleRef{罗 1:25}我想，你会把洋葱和块菌也算作地的恩惠，因为主宣布：\BibleQuote{人生活不只靠饼！}\BibleRef{玛 4:4}这样，犹太人因将希望局限于地上的事物而失去了天上的福乐，他们不知道天上之饼的许诺，不知道天主傅油之油、圣神之酒，以及那从基督之葡萄树获得活力的生命之水。基于完全相同的原则，他们视\Place{犹大}的独特土壤为那圣洁之地，而这圣洁之地更应解释为主肉身，即在所有穿上基督的人身上，此后就成为圣洁之地；确实因圣神的居留而圣洁，因祂保证的甘美而真正流奶流蜜，因天主的友谊而真正是犹大。因为\BibleQuote{外表上的犹太人不算犹太人，只有内心的才是犹太人。}\BibleRef{罗 2:28}同样，当\Person{依撒意亚}说：\Quote{醒来，醒来，耶路撒冷！穿上你手臂的力量；醒来，如同你最初之时}——即指在犯罪堕落之前的无罪状态，这天主的圣殿和耶路撒冷也必须以此方式理解。因为这种劝勉和邀请的话，怎能适合于那杀害先知、用石头砸死被派遣者、最后甚至钉死自己主的耶路撒冷呢？救恩也确实未曾许诺给任何一片土地，因为土地必随着整个世界的形态而消逝。即使有人大胆主张，乐园就是那圣洁之地——或许可称为我们原祖\Person{亚当}和\Person{厄娃}之地——那么即使如此，乐园的恢复也似乎是对肉身的许诺，因为肉身正是受命居住并看守乐园的，为使人能被召回那里，正如他被逐出时那样。\par

% structure: p0055 source=h2 target=section reason=relative-heading-rank; base=h2

\section{第二十七章 某些比喻性术语就肉身复活加以解释}

我们在圣经中也看到\Emph{长袍}被提及，作为肉身希望的预像。这样在\BibleBook{若望}{默示录}中说：\Quote{这些是没有用女人污秽自己衣服的人}——当然这是指童贞者，以及那些成为\Quote{为天国自阉的人}。\BibleRef{玛 19:12}因此他们将\Quote{穿上白衣}，\BibleRef{默 3:5}即未嫁娶之肉身的明亮美丽。在福音中，甚至\Quote{婚宴礼服}也可被视为肉身的神圣。\BibleRef{玛 22:11}因此，当依撒意亚告诉我们主所拣选的何种斋戒，并加上关于善行赏报的陈述时，他说：\Quote{那时你的光明要如晨光出现，你的衣服要迅速升起}；\BibleRef{依 58:8}这里他并非想到外衣或布袍，而是指肉身的升起，即他在死亡堕落之后所宣告的复活。这样我们甚至得到了身体复活的预像性辩护。因此，当我们读到：\Quote{我的百姓，进入你的密室，停留片刻，直到我的怒气过去}，\BibleRef{依 26:20}我们就把密室理解为坟墓，那些在世界的终末、敌基督势力的最后猛烈进攻中离开此世的人将在此安息片刻。不然他为何使用\Emph{密室}一词而不用其他容器，难道不是因为这肉身被盐腌在这些密室或地窖中保留备用，以便在合适的时候从那里取出吗？同理，涂香料的尸体也被存放在陵墓和坟墓中待葬，以便当主人命令时可以从那里移出。因此，既然如此理解经文有一致性（因为小小的密室如何能庇护我们免受天主的义怒？），\Emph{看来}正是借着所使用的词句\Quote{直到他的怒气过去}——那将消灭敌基督的，他事实上表明，在那愤怒之后，肉身将从坟墓中出来，这肉身在怒气\Emph{爆发}之前就已存放在那里。如今从密室中取出的无非是放入其中的东西，而且在敌基督被铲除之后，将忙碌地施行复活\Emph{的伟大过程}。\par

% structure: p0057 source=h2 target=section reason=relative-heading-rank; base=h2

\section{第二十八章　先知的事物和行动，以及言语，都证明这伟大的教义。}

但我们知道，预言不仅藉言语表达，也藉事物表达。复活既由言语，也由行为预告。当梅瑟将手放入怀中，又抽出来，手已死了；再放入怀中，又抽出来，手复活着，\BibleRef{出 4:6}这不就是全人类的预兆吗？——因为那三个标记\BibleRef{出 4:2}表示天主的三重能力：首先，按预定顺序，使人类征服那凶猛的古蛇、魔鬼；其次，将肉身从死亡的怀中抽出；最后，在审判中追讨所有流出的血。关于这主题，我们在同一先知的著作中读到，天主说：\Quote{我要向一切野兽追讨你们生命的血；也要向人手和他兄弟的手追讨。}\BibleRef{创 9:5}现在，什么也不被要求，除非是那被要求归还的；而所要求的无非是那必须交出的；那交出之物，当然就是基于报复而被要求和索还的。但实际上，对从未存在过的事物不可能有惩罚。然而，当它被恢复以便受惩罚时，它将拥有存在。因此，凡关于血所宣告的一切，都适用于肉身，因为没有肉身就没有血。肉身将被复活，以便血可以受惩罚。还有一些圣经陈述如此明确，毫无寓言的晦涩，然而它们强烈要求按其直白来理解。例如，依撒意亚中的那段话：\Quote{我杀死，也使人活。}当然，他使人活是在他杀死之后。因此，正如他是藉死亡来杀害，他也是藉复活来使人活。现在是肉身被死亡杀害；因此，肉身将被复活所复苏。的确，如果杀害意味着从肉身夺去生命，而其对立的复苏等于将生命恢复给肉身，那么肉身必定要复活，因为被杀害夺去的生命必须通过复活来恢复。\par

% structure: p0059 source=h2 target=section reason=relative-heading-rank; base=h2

\section{第二十九章　引用厄则克耳关于枯骨的异象。}

因此，既然圣经的比喻部分、事实的论证、以及某些\Emph{圣经}的明确陈述，都光照了肉身的复活（虽然并未特别提到那实体本身），那么，那些明确提及身体实体本身的经文，岂不更有效地确定这个问题吗！以\Person{厄则克耳}为例：\BibleQuote{主的手，}他说，\BibleQuote{就放在我身上；主在圣神内带我出去，把我放在一个布满骨头的平原中间；祂领我绕着他们走了一圈：看，平原的表面上有很多骨头；看，它们非常枯干。祂对我说：人子，这些骨头能活吗？我回答说：上主天主，祢知道。祂对我说：你向这些骨头讲预言，对他们说：枯干的骨头，听上主的话！上主天主对这些骨头这样说：看，我要使\Emph{生命}的气息进入你们内，你们必要复活；我要赐给你们气息，在你们身上安上肌肉，包上皮肤；你们必要复活，并且知道我是上主。我就照主所命令的预言了；正预言时，忽然有响声，有震动，骨头互相靠近。我看见，看，筋和肉长在上面，肌肉包着它们；但里面没有气息。祂对我说：人子，你向风讲预言，说：上主天主这样说：气息啊，你要从四方吹来，吹在这些被杀的人身上，使他们复活。我就像祂命令我的那样向风预言，气息就进入骨头，他们就复活了，站起身来，成为一支极其庞大的军队。\Emph{主}对我说：人子，这些骨头就是以色列全家。他们自己说：我们的骨头枯干了，我们的希望消失了，我们在这骨头中被彻底毁灭了。因此你向他们讲预言，说：看，我要亲自打开你们的坟墓，把你们从坟墓中领出来，我的子民，领你们进入以色列地；你们就知道我是上主，我开了你们的坟墓，把你们从坟墓中领出来，我的子民；我要把我的神赐给你们，你们必要复活，在你们自己的土地上安居；你们就知道我上主说了并行了这些事，这是上主说的。}\BibleRef{则 37:1}\par

% structure: p0061 source=h2 target=section reason=relative-heading-rank; base=h2

\section{第三十章 \Person{戴尔都良}对此异象的解释：论死者身体的复活。\newline{}本作者的一个年代错误：他以为\Person{厄则克耳}在其第三十一章中所预言的是在被掳之前。}

我深知他们如何甚至将这一预言扭曲为寓意的证据，理由是当主说：\BibleQuote{这些骨头是以色列全家}时，祂使之成为以色列的象征，并使之脱离其本来的字面状况；因此（他们争辩说）这里是对复活的一种比喻性预言，而非真实的预言，因为（他们说）犹太人的状况是某种意义上的卑微、枯干，且分散在世界平原上。因此复活的形象被寓意式地应用于他们的状况，因为他们需要被聚集、骨与骨相连（换句话说，支派与支派相联、民族与民族相合），并通过权力的筋和君王的神经得以重新结合，并从坟墓中出来（即从最悲惨、最卑贱的俘囚居所中出来），在复兴的道路上重新呼吸，并从此生活在自己的犹大地。这一切之后会发生什么？他们无疑会死去。死后又会怎样？当然没有从死者中的复活，因为这里向厄则克耳启示的没有这类事。然而，复活在别处已被预言：所以即使在这种情况下也将有一次复活，而他们竟敢将\Emph{这段经文}应用于犹太事务的状况是鲁莽的；或者，即使它指示了一种与我们主张的复活不同的复原，那与我何干，只要也有身体的复活，正如犹太国度的复兴一样？事实上，正是由于犹太国度的复兴是由骨头的重新结合与接连所预表，这就证明了这事也将发生在\Emph{骨头本身}；因为若不是同样的事确实要在骨头中实现，就不能从骨头形成隐喻。然而，尽管在形象中有了真实事物的草图，形象本身仍拥有自身的真理：因此，被借用来比喻其他事物的事物，必须先有自身的存在。空虚不能作为比喻的一贯基础，无有也不能作为比喻的合适根基。因此，应当相信骨头注定要重新获得肉和气息，正如\Emph{这里}所说的那样，正因如此，它们更新的状态才能单独表达犹太事务的改革状况——而这被认为是这段经文的含义。不过，更符合虔诚精神的做法是，根据字面解释的权威来维护真理，正如圣神所默感的经文的意义所要求的那样。现在，如果这神视关乎犹太人的状况，那么主在向他启示骨头的位置之后，会立即加上：\BibleQuote{这些骨头是以色列全家}，等等。但是，在展示骨头之后，祂立即中断场景，说了一些最适合骨头的前景；尚未提及以色列时，祂试探先知本人的信德：\BibleQuote{人子，这些骨头能复活吗？}以便他回答：\BibleQuote{主，祢知道。}你可知，天主绝不会在一个永远不真实、以色列永远不会听到、\Emph{也}不应相信的点上试探先知的信德。然而，既然死者复活的确已被预言，但以色列因其极大不信的疑虑而对此反感；并且，当他们凝视朽坏坟墓的状况时，对复活绝望；或者更确切地说，他们并没有主要将心思放在这上面，而是放在他们自己困扰的处境上——因此天主首先教导了先知（因为他自己也并非没有疑虑），通过向他启示复活的过程，以便他热切地宣扬同样的道理。然后祂命令百姓相信祂向先知启示的事，告诉他们，他们自己——尽管拒绝相信他们的复活——正是注定要复活的那些骨头。最后，在结尾句中祂说：\BibleQuote{你们要知道，我上主说了并行了这些事}，当然用意是要做祂所说的事；但如果祂的设计是要做与祂所说的不同的事，那么祂决不是打算做祂所说的事。\par

% structure: p0063 source=h2 target=section reason=relative-heading-rank; base=h2

\section{第三十一章 先知书中其他关于肉身复活的经文。}

毫无疑问，如果人民沉溺于\Emph{比喻的}怨言，说他们的骨头枯干，他们的希望破灭——因他们分散的结果而哀叹——那么天主用\Emph{比喻的}应许来安慰他们的\Emph{比喻的}绝望，似乎相当合理。然而，既然人民尚未因分散而遭受任何伤害，尽管复活的希望在他们中间常常落空，显然是因为他们身体的朽坏状态动摇了他们对复活的信德。因此，天主正在重建人民所拆毁的信德。但即使以色列当时因现有处境的一些打击而沮丧，我们也不应因此认为启示的目的可以停留在比喻中：它的目的一定是证实复活，以便将民族的希望提升到永生的救恩和不可或缺的恢复，从而转移他们对当前事务的忧虑。这确实也是其他先知的目的。\BibleQuote{你们要出去，}（\Person{玛拉基亚}说）\BibleQuote{从你们的坟墓出来，如同挣脱绳索的牛犊，你们要践踏你们的仇敌。}\BibleRef{拉 4:2} 又一次，（\Person{依撒意亚}说）：\BibleQuote{你们的心要喜乐，你们的骨头要像草一样发芽，}\BibleRef{依 66:14} 因为草也是通过种子的分解和腐朽而更新的。总之，如果有人认为复活骨头的比喻直接指向以色列的状况，为什么同样的希望宣布给万民，而不限于以色列，即重新赋予那些骸骨以身体和生命的气息，并从坟墓中复活他们的死人？因为这话是普遍的：\BibleQuote{死人要复活，从坟墓中出来；因为从你而来的露水是她们骨头的良药。}\BibleRef{依 26:19} 在另一处\Emph{记载着}：\BibleQuote{凡有血肉的，都要来在我面前敬拜，主说。}\BibleRef{依 66:23} 何时呢？当这世界的样式开始消逝时。因为祂先前说过：\BibleQuote{如同我所造的新天新地在我面前存留，主说，你们的后裔也要如此存留。} 那时也要应验随后所写的：\BibleQuote{他们要出去}（即从他们的坟墓），\BibleQuote{并且看见那些悖逆之人的尸首：因为他们的虫子不死，他们的火也不灭；他们将成为一切血肉的景象}\BibleRef{依 66:24} 甚至包括那些从死人中复活、从坟墓出来、为这宏大的恩宠而朝拜主的人。\par

% structure: p0065 source=h2 target=section reason=relative-heading-rank; base=h2

\section{第32章 即使未埋葬的尸体也将复活。无论遭遇什么，天主都将恢复它们。\Person{约纳}的案例被引用来说明天主的权能。}

但是，为使你不认为只有那些被安葬在坟墓中的尸体才被预言复活，你在圣经中读到：\BibleQuote{我要命令海中的鱼，它们必吐出所吞吃的骨头；我要使关节对关节，骨头对骨头。} 你可能会问：那么鱼和其他动物以及食肉的鸟也会复活，以便吐出它们所吞吃的，因为你在\Person{梅瑟}法律中读到，甚至要向所有野兽追讨血债？当然不是。但提到野兽和鱼是与血肉的恢复有关，为了更加强调即使被吞吃的身体也会复活，据说要向它们的吞吃者追讨补偿。现在我认为在\Person{约纳}的事例中，我们有天主的权能的公平证明，当他从鱼腹中出来，他的两个本性——他的肉体和灵魂——均未受损伤。毫无疑问，鲸鱼的内脏在三天内有充足时间消耗和消化\Emph{约纳的}肉体，完全如同棺材、坟墓或某个安静隐蔽坟墓的逐渐腐朽；只是他想预表那些野兽（象征）特别强烈反对\Emph{基督徒的}名字的人，或邪恶天使，对于他们，复仇审判将完全追讨血债。那么，哪里有这样一个人，他更倾向于学习而非假设，更谨慎相信而非争论，对天主的智慧更谨慎而非任性固执于自己的智慧，当他听到有关筋腱、皮肤、神经和骨头的天主的计划时，会立刻想出对这些话的不同应用，仿佛所有这些关于所讨论实体的说法都不是自然指向人类？因为要么这里没有指向人的命运——在天国恩慈的预备中，在审判之日的严厉中，在复活的所有事件中；要么，如果指向他的命运，那么目的地必须参照构成人的那些实体，命运是为这人保留的。我还有一个问题要问这些非常灵巧的转化骨头、筋腱、神经和坟墓的人：为什么当宣告关于\Emph{灵魂}的任何事时，他们不把灵魂解释为别的东西，并将其转移到另一个含义？——因为每当对\Emph{身体的}实体作出明确陈述时，他们顽固地喜欢采用任何其他意义，而不是名称所指的意义。如果属于身体的事物是比喻的，为什么属于灵魂的事物不是比喻的呢？然而，既然属于灵魂的事物没有任何寓意，那么属于身体的事物也没有。因为人是身体正如他是灵魂；所以不可能这两种本性之一接受比喻的意义，而另一个排除它。\par

% structure: p0067 source=h2 target=section reason=relative-heading-rank; base=h2

\section{第33章 关于先知书就说这么多。在福音中，\Person{基督}亲自解释的比喻清楚地指向肉体的复活。}

由先知书，这已是足够的证据。现在，我诉诸于福音书。但在此，我也必须首先应对同样的诡辩，这诡辩是由那些主张主如同（先知）一样、以寓意方式说话的人提出的，因为经上记载：\BibleQuote{耶稣用比喻对他们讲论这一切，不用比喻就不给他们讲什么。}\BibleRef{玛 13:34}，即对犹太人。门徒也问祂：\BibleQuote{祢为什么用比喻对他们讲话？} 主回答他们说：\BibleQuote{我之所以用比喻对他们讲话，是因为他们看，却看不见；听，却听不见，正应验了依撒意亚的预言。} 但因为祂是用比喻对犹太人讲话，那就不是对所有人；如果不是对所有人，那就意味着祂并非总是在所有事情上用比喻，而只是在某些事情上，且在针对某一特定群体时。但当祂对犹太人讲话时，祂是针对一个特定群体。诚然，祂有时甚至对门徒也用比喻。但请注意圣经如何记载这类事实：\BibleQuote{祂对他们讲了一个比喻。} 由此可见，祂并非通常用比喻对他们讲话；因为如果祂总是这样做，就不会特别提到祂采用这种讲话方式。此外，没有一个比喻是你不能发现要么由主亲自解释的，如同撒种的比喻（祂解释为关于天主圣言的管理）；要么由福音作者在前言中阐明，如同关于傲慢的法官和恳切的寡妇的比喻，\Emph{明确地被应用于}祈祷的恳切；要么能够自然地被理解，如同关于无花果树的比喻，它被暂时存留，希望有所改善——这是犹太人不结果实的象征。如今，如果连比喻也不遮蔽福音的光明，那么那些含义明确、无可置疑的句子和声明，怎么可能表示与字面意义不同的东西呢！但主正是通过这样的声明和句子来陈明末后的审判、天主的国或复活的：\BibleQuote{在审判之日，对于提洛和漆冬，}祂说，\BibleQuote{\Emph{比你们}更可忍受。}\BibleRef{玛 11:22} 和\BibleQuote{告诉他们：天主的国临近了。}\BibleRef{玛 10:7} 并且，\BibleQuote{在义人复活时，你必得报答。}\BibleRef{路 14:14} 如今，如果这些事件（我指审判之日、天主的国和复活）的提及有明确而绝对的意义，以至于关于它们没有任何东西能被牵强地解释为寓意，那么那些描述\Emph{天主的}国的安排、过程和体验，以及审判和复活的陈述，也不应被强行解释为比喻。相反，那些注定属于身体的事物应当谨慎地按身体的意义来理解——而不是按属神的意义，因为它们本质上没有任何比喻性。这就是为什么我们预先设定了一个初步考虑，即灵魂和肉身两者的身体实体都可能得到报应，这报应将根据两种本性的合作而判给，以便灵魂的形体性不会通过暗示求助于比喻性描述而排除肉身的身体性质，因为两者都必须被视为注定要参与天主的国、审判和复活。现在，我们继续对此命题的专门证明：每当我们的主提及复活时，祂都指出了肉身的身体特征，同时并未贬低灵魂的形体本质——这一点实际上只有少数人承认。\par

% structure: p0069 source=h2 target=section reason=relative-heading-rank; base=h2

\section{第34章 基督明确为整个人的复活作证。不只是在灵魂，而没有身体。}

先来看祂说祂来是\BibleQuote{\Emph{寻找并}拯救那喪亡的}的那段话。\BibleRef{路 19:10}你以为那喪亡的是什么？无疑是人。是整个的人，还是他的一部分？当然是整个人。事实上，由于导致人喪亡的过犯，既由灵魂因私欲偏情的煽动，也由肉身因实际享乐的行动而犯下，它已将整个的人置于过犯的判决之下，因此理当使他沦于喪亡。所以他将被完全拯救，因为他因犯罪而完全喪亡了。除非（比喻中的）羊是一只\Quote{喪亡的}羊，而与它的身体无关；那么找回它可能不需要身体。然而，既然构成全羊的身体和灵魂，都被善牧扛在肩上，我们这里无疑就有了一个例子，说明人是如何在其两种本性中恢复的。否则，天主只带人一半（甚至更少）到救恩中，这多么有损于祂的尊严——而现世君王的慷慨总是以完全的恩宠自居！那么，魔鬼伤害人、使人完全喪亡，就必须被认为更强有力吗？而天主则显得相对软弱，因为祂不帮助人于其完整状态？然而宗徒提示说：\BibleQuote{罪恶在哪里增多，恩宠在那里也格外丰富。}\BibleRef{罗 5:20}事实上，一个人若同时被说成是喪亡的——即肉身喪亡，而灵魂得救，他怎能被视为得救呢？除非\Emph{他们的论辩}现在要求灵魂被置于\Quote{喪亡}状态，以便它可以接受救恩，理由是只有喪亡的才真正得救。然而我们理解灵魂的不朽，是相信它\Quote{喪亡}，不是指毁灭，而是指惩罚，即在地狱里。若是如此，那么救恩将不涉及灵魂，因为它已因其本性不朽而\Quote{安全}了；而涉及肉身，这肉身，正如众人轻易承认的，是会朽坏的。否则，如果灵魂也是可朽的（在此意义上），即不死的——与肉身的情况相同——那么这相同的情况按理也应当有利于肉身，因为它同样是可朽和可喪亡的，因为主打算拯救那喪亡的。我现在不想继续我们讨论的线索，去考虑丧亡是在人的一种本性还是在另一种本性中提出它的要求，只要救恩均匀地分布在两种实体上，并以两者为目标。因为请注意：你假设人在哪一种实体中喪亡了，在另一种实体中他就不喪亡。因此他将在他不喪亡的实体中得救，却又在他喪亡的实体中获得救恩。你便有了全人的恢复，因为主打算拯救他喪亡的那部分，而当然不会失去那不能喪亡的部分。谁还会怀疑两个本性的安全呢？其中一个获得救恩，另一个则不失去它。此外，主在说以下的话时向我们解释了这件事的意义：\BibleQuote{我而来不是为执行我的旨意，而是为执行派遣我来的父的旨意。}\BibleRef{若 6:38}我问，那旨意是什么？\Quote{凡祂赐给我的，叫我连一个也不失掉，反而在末日使他复活。}现在，当祂更强调地重复祂的话时：\Quote{凡看见子，并信从子的，必获得永生；并且在末日，我要使他复活。}——祂肯定了复活的全备性。因为祂将适合其服务的报酬分派给每一个别的本性：给肉身，因为子被它\Quote{看见}；给灵魂，因为子被它\Quote{信从}。那么，你会说，这应许是给那些基督被他们\Quote{看见}的人。好吧，就如此；只愿同样的希望从他们流到我们！因为如果对那些看见因而相信的人，如此果实从肉身和灵魂的操作中产生，那么对我们更是何等大呢！因为基督说：\BibleQuote{那些没有看见而相信的，才是有福的。}\BibleRef{若 20:29}既然，即使肉身的复活必须对\Emph{他们}否认，它也至少应当是\Emph{我们这些更为有福的}合宜恩赐。因为我们怎能有福，如果我们任何部分丧亡了呢？\par

% structure: p0071 source=h2 target=section reason=relative-heading-rank; base=h2

\section{第三十五章 论何谓将复活的身体。并非灵魂的形体性。}

但祂也教导我们：\Quote{\InnerQuote{那能在地狱中毁灭身体和灵魂的，更当畏惧祂}}，就是说，唯有主自己；\Quote{\InnerQuote{那些只杀害身体而不能伤害灵魂的}}（\BibleRef{玛 10:28}），即指所有属人的权势。如此，我们便承认灵魂本性不死，人不能杀害它；而身体是会死的，人可能杀害它：\Emph{我们由此得知}死人复活乃是肉身的复活；因为若不再复活，肉身就不可能\Quote{\InnerQuote{在地狱中被杀}}。但可能有人会吹毛求疵地质问\Quote{\InnerQuote{身体}}（或\Quote{\InnerQuote{肉身}}）的含义，我当即声明：我所谓的身体，无非就是那由血肉构成，不论其构造和改变的材料为何，为人所见、所触摸，有时甚至被人杀害的东西。同样，我不会承认墙的\Emph{身体}除了灰泥、石头和砖块之外还有什么。若有人向我们的论证引进某种精微隐秘性质的身体，他必须向我表明、揭示、证明那同一身体正是被人间暴力所杀害的身体，然后（我才同意）我们的经文所说的是这样的身体。再者，若有人拿灵魂的身体\Emph{或形体性}来反驳我，那不过是无聊的藉口！因为既然（在这一段经文中）两种实体都呈现在我们面前，说\Quote{身体和灵魂}都在地狱中被毁灭，那么两者之间显然有区别；我们只能认为身体是我们可触摸的，即肉身，它既然没有\Quote{更畏惧}被天主毁灭，因而在地狱中被毁灭，同样也将恢复永生，因为它宁愿被人手杀害。所以，若有人强行认为灵魂和肉身在地狱中的毁灭意味着两种实体的最终湮灭，而非其受罚的待遇（仿佛它们是被耗尽，而非受惩罚），那么请他记起地狱之火是永恒的——明确宣布为永恒的惩罚；然后请他承认，正是因此，这无尽的\Quote{杀戮}比仅仅暂时的世人的杀害更为可怕。他便会得出结论：实体必须是永恒的，因为对其惩罚性的\Quote{杀戮}是永恒的。既然复活后的身体必须与灵魂一起在地狱中被天主所杀，我们显然从这一事实中获得充分的信息，关于身体将要面对的\Emph{两方面结局}，即肉身的复活及其永恒的\Quote{杀戮}。否则，若肉身被复活而注定要\Quote{\InnerQuote{在地狱中被杀}}以便被终结，那将是极其荒谬的，因为它若不复活，本可（更直接地）承受这种湮灭。这真是奇怪的悖论：为了让一种本质获得实际上已经加给它的湮灭，还必须重新赋予它生命！\Emph{但基督}在坚固我们同一盼望时，加上了\Quote{\InnerQuote{麻雀}}的例子——\Quote{\InnerQuote{一只也不会不带天主的旨意掉在地上}}（\BibleRef{玛 10:29}）。\Emph{祂说这话}，是叫你相信，同样地，被弃于地上的肉身能因同一天主的旨意而复活。因为虽然这对麻雀而言并不允许，但我们\Quote{\InnerQuote{却比许多麻雀贵重}}，恰恰因为当我们跌倒时，我们复活了。最后，祂肯定说：\Quote{\InnerQuote{连我们的头发也都数过了}}（\BibleRef{玛 10:30}），在肯定中当然也包含保全它们的许诺；因为若它们都将失落，如此仔细地数点它们又有什么用？当然，唯一用处在于这一真理：\Quote{\InnerQuote{凡父赐给我的，我必不失落任何一个}}（\BibleRef{若 6:39}）——连一根头发也不失落，同样也不失落一只眼或一颗牙。然而，那\Quote{\InnerQuote{哀哭和切齿}}将从何处来呢？如果不是来自\Emph{眼睛和牙齿}？——正是在那时，身体在地狱中被杀，被投入外面的黑暗中，那黑暗将是眼睛合适的折磨。那在婚宴上未穿着善行之衣的人，将被\Quote{\InnerQuote{捆绑手脚}}——当然，是在他的身体复活之后。同样，在天主的国中坐席、坐在基督的宝座上、最后站在祂左右、吃生命树上的果子：这些若不是身体性安排和归宿的最确实证据，又是什么呢？\par

% structure: p0073 source=h2 target=section reason=relative-heading-rank; base=h2

\section{第三十六章 基督驳斥撒杜塞人，肯定天主教教义。}

现在让我们看看，主是否在破除撒杜塞人狡猾的诡辩时，没有赐予我们的教义更大的力量。我认为，他们的主要目标是完全取消复活，因为撒杜塞人事实上既不承认灵魂的救恩，也不承认肉身的救恩；因此，他们采用最能削弱复活可信度的案例，从中提出一个论点来支持他们所提出的问题。他们似是而非的问题涉及肉身，即复活后是否还要受婚姻约束；他们假设了一个嫁过七个兄弟的妇女的情况，以至于她应该归于哪一个，成了一个疑点。现在，让我们牢牢把握问题和答案的要旨，讨论便立即解决了。因为撒杜塞人确实否认复活，而主却肯定复活；同样，（在肯定时，）祂指责他们既不知道圣经——当然是指那些宣告复活的圣经——也不相信天主的大能，尽管天主的大能当然足以使死人复活；最后，因为祂紧接着又毫无疑虑地说：\BibleQuote{至于死者复活}\BibleRef{路 20:37}，并肯定那正被否认的事，即在祂这位\Quote{活人的天主}面前死者的复活——（显然可以得出）祂正是在他们否认的确切意义上肯定了这一真理；事实上，这就是人的两种本性的复活。也不像他们所以为的那样，因为基督否认人将结婚，祂就证明了人不会复活。相反，祂称他们为\Quote{复活之子}，在某种意义上，他们借着复活要经历一次出生；之后他们不再结婚，在复活的生命中\Quote{与天使平等}，因为他们不再结婚，也因为他们不再死亡，而是注定要穿上不朽的衣服进入天使状态，尽管复活的实体有所改变。此外，任何关于我们是否再次结婚或死亡的问题，都不可避免地会牵扯到对那与死亡和婚姻都有特殊关系的实体——即肉身——的恢复的怀疑。因此，你们看到主反对犹太异端者所肯定的，正是如今基督徒中的撒杜塞人所否认的——整个人的复活。\par

% structure: p0075 source=h2 target=section reason=relative-heading-rank; base=h2

\section{第三十七章 基督关于肉身无益的断言与我们教义的一致性解释。}

他说，的确，\BibleQuote{肉一无所用；}\BibleRef{若 6:63} 但随后，如在前例中，意义须按所谈论的主题来定。现在，因为他们认为他的话生硬难堪，以为他真的命令他们吃他的肉，为将救恩状态安排为属神的事，他先提出原则：\BibleQuote{使人生活的是神；}然后加上：\BibleQuote{肉一无所用，}——当然是指赐予生命而言。他继续解释他愿意我们如何理解\Emph{神}：\BibleQuote{我对你们所说的话，就是神，就是生命。}类似地，他先前曾说：\BibleQuote{听我的话，相信派遣我来的那一位，就有永生，不受审判，而是已出死入生了。}\BibleRef{若 5:24} 因此，他将自己的话立为赐生命的原则，因为那话是神和生命，他也同样用这名称呼他的肉；又因为圣言成了血肉，\BibleRef{若 1:14} 所以我们应当渴慕他，好得到生命，用耳朵吞噬他，用理解力反复思想，用信德消化他。就在所谈论的经文之前，他曾宣布他的肉是\BibleQuote{从天上降下来的食粮，}\BibleRef{若 6:51} 不断用必要食物的比喻，在听者心中铭刻他们祖先的回忆，他们曾偏爱埃及的食粮和肉胜过天主的召唤。然后，他将话题转向他们的反思，因为他觉察到他们将从他那里散开，他说：\BibleQuote{肉一无所用。}现在，有什么能摧毁肉身复活呢？仿佛不可能合理存在某物，虽然它本身\Quote{一无所用}，却可能被另一物所用。神\Quote{有益}，因为神赐予生命。肉一无所用，因为肉属于死亡。因此，他反而以有利于我们信仰的方式提出了两个命题：因为通过显示什么\Quote{有益}，什么\Quote{无益}，他也阐明了接受者和给予者两方。因此，\Emph{在此情形中，我们有}神赐生命给被死亡征服的肉身；因为\Quote{时辰}，他说，\BibleQuote{到了，那时死人都要听见天主子的声音，凡听见的都要生活。}\BibleRef{若 5:25} 现在，\Quote{死人}不是肉身吗？\Quote{天主的声音}不是圣言吗？圣言不是神吗？那神将公义地复活他曾亲自成为的肉身，而且\Emph{那}肉身也脱离死亡（他自己曾受的死亡）和坟墓（他亲自进入的坟墓）？再者，当他说：\BibleQuote{你们不要惊奇这事，因为时候要到，那时凡在坟墓里的都要听见天主子的声音，并且出来；行过善的，复活进入生命；作过恶的，复活接受审判。}\BibleRef{若 5:28}——在这些话之后，无人能将\BibleQuote{在坟墓里的}死人解释为别的，只能是肉身之躯，因为坟墓本身无非是尸体的安息之所：因为无可争辩的是，甚至那些有份于\Quote{旧人}的，即罪人——换句话说，那些因不认识天主而死亡的人（我们的异端者竟然愚蠢地坚持用\Quote{坟墓}一词来理解他们）——在这里明确被说成要从他们的坟墓中出来受审判。但坟墓如何从坟墓中出来呢？\par

% structure: p0077 source=h2 target=section reason=relative-heading-rank; base=h2

\section{第三十八章 基督藉复活死人，以实际行动证实肉身复活的道理。}

在主的\Emph{言语}之后，对于祂\Emph{行动}的意图，当祂从棺架和坟墓中复活死人时，我们应如何思考？祂为何这样做？如果只是为了展示祂的能力，或给予暂时恢复生命的恩惠，那么复活人后又让他们再次死亡，实在不是什么大事。然而，如果正如事实所是，祂更是在稳妥地保存人们对未来复活的信仰，那么根据祂自己例证的特定形式，必然得出所说的复活将是身体的复活。我绝不允许有人说，未来的复活，既然只为灵魂而设，那时却先得到了这些肉体复起的初步例证，仅仅是因为除非通过可见实体的复活，否则无法显示不可见灵魂的复活。那些认为天主只能做他们思想范围内之事的人，对天主的认识实在贫乏；毕竟，只要他们熟悉\Person{若望}的著作，就不能不全知祂的能力曾经如何。因为毫无疑问，那位向我们展示了殉道者迄今无体的灵魂安息在祭台之下的（\BibleRef{默 6:9}），完全有能力让我们看到他们没有肉身的身体而升起。然而，我宁愿（相信）天主不可能行欺骗（祂只有在诡计方面才是软弱的），因为害怕似乎以一种与祂实际处置事物不一致的方式给予事物初步的证明；更进一步，因为害怕，既然祂没有足够的能力向我们展示无肉体的复活样本，祂可能更加软弱无力地无法在不久后于同一\Emph{肉体}实体中展示样本的完全实现。的确，没有例证比它作为样本的事物更大。然而，如果灵魂与其身体同被复起，作为它们无体复活的证据，以至于人的整个救恩\Emph{在灵魂和身体上}变成仅一半的保证，即\Emph{灵魂}，那例证就更大了；而所有例证的情况是，被认为较小者——我指的是仅灵魂的复活——应作为肉体在指定时间复起的预尝。因此，按照我们对真理的估计，那些被主复活的死人的例证确实是肉体与灵魂复活的证明——事实上，证明这恩赐不会被拒绝给任何一种实体。然而，仅作为例证考虑，它们表达出的意义要小得多——\Emph{确实更小}，比基督最终要表达的小——因为它们不是为荣耀和不朽而复活，而是只为另一次死亡。\par

% structure: p0079 source=h2 target=section reason=relative-heading-rank; base=h2

\section{第三十九章 \WorkTitle{宗徒大事录}给我们提供的额外证据}

\WorkTitle{宗徒大事录}也证实了复活。宗徒们在犹太人中间，除了解释旧约并确认新约，以及首要的是在基督内宣讲天主之外，别无他事。因此，关于复活，除了将其宣告归于基督的光荣外，他们并没有引入什么新东西；在其它方面，它早已被单纯而明智的信仰所接受，对于复活将是何种方式没有任何疑问，除了撒杜塞人之外也没有遇到任何反对者。完全否认复活比以一种异样的方式理解它要容易得多。你会发现\Person{保禄}在千夫长的庇护下，在司祭长、撒杜塞人和法利塞人面前宣认自己的信仰：\Quote{诸位仁人弟兄，}他说，\Quote{我是法利塞人，是法利塞人的儿子；我现今受审，是因死人的望德和死人复活，}（\BibleRef{宗 23:6}）——当然是指民族的望德；为了避免在他当时的情况下，作为一个明显的违法者，被认为在信仰最重要的一点上——即复活——接近撒杜塞人的意见。因此，他表面上不愿损害对复活的信仰，实际上却是确认了法利塞人的意见，因为他拒绝了否认复活的撒杜塞人的观点。同样，在\Person{阿格黎帕}面前，他也说他所宣讲的\Quote{没有别的，就是先知们所预告的。}（\BibleRef{宗 26:22}）因此，他所维护的正是先知们所预言的复活。他还提到\Quote{梅瑟}所写的关于死人复活的事（并这样做时），他必定知道那将是身体的复活，因为那里要对人的血进行追究。（\BibleRef{创 9:5}）他于是宣称，这复活具有法利塞人所承认的性质，也是主自己所维护的性质，并且是撒杜塞人拒绝相信的性质——这种拒绝导致他们实际上完全否认真理。先前，雅典人也没有理解\Person{保禄}宣布任何其他的复活。（\BibleRef{宗 17:32}）事实上，他们曾嘲笑他的宣告；但如果他们从他那里只听到灵魂的恢复，他们就不会如此嘲笑，因为他们会把\Emph{那}视为他们本土哲学非常普通的预期。但当关于复活的宣讲——他们先前未曾听过——以其绝对的新颖激动了外邦人，一种并非不自然的对如此奇妙之事的不信开始以许多讨论困扰单纯的信仰时，宗徒便在他几乎每一部著作中都注意加强人们对这\Emph{基督徒}望德的确信，指出有这种望德，它尚未实现，并且将是在身体内——这是特别研究的对象，此外还有一个疑难问题，即不是在与我们不同的身体内。\par

% structure: p0081 source=h2 target=section reason=relative-heading-rank; base=h2

\section{第四十章 圣保禄多处经文证明我们的教义，从异端的歪曲中抢救出来}

如今，如果有人巧妙地从（宗徒）的著作中挑出论证来，这并不奇怪，因为\BibleQuote{必须有异端；}\BibleRef{格前 11:19}但如果圣经不可能被错误解释，这些就不可能发生。那么，异端者发现宗徒提到了两个\Quote{人}——\Quote{内里的人，}就是灵魂，以及\Quote{外面的人，}就是肉身——便将救恩归于灵魂或内里的人，将毁灭归于肉身或外面的人，因为（在致格林多人的书信中）记载：\BibleQuote{我们外面的人虽然朽坏，但我们内里的人却日日更新。}\BibleRef{格后 4:16}如今，仅灵魂本身并不是\Quote{人}（它是后来被植入那已被称为\Emph{人}的泥土模型中的），没有灵魂的肉身也不是\Quote{人}：因为灵魂离开后，它被称为尸体。因此，\Emph{人}这一称谓在某种意义上是两种紧密结合实体之间的纽带，在这一称谓下，它们只能是相连的本性。至于内里的人，宗徒更愿意将其视为心思和心，而非灵魂；换句话说，不是实体本身，而是实体的味道。因此，当他致书厄弗所人时，说到\BibleQuote{基督住在他们内里的人中，}\BibleRef{弗 3:17}无疑他意指主应当被接纳到他们的感官中。然后他接着说：\BibleQuote{在你们心中藉着信德，在爱德中\Emph{根深蒂固}，}——将\Quote{信德}和\Quote{爱德}视为不是实体的部分，而只是灵魂的意念。但当他说到\Quote{在你们心中}时，由于心是肉身的实体部分，他立刻将实际的\Quote{内里的人}归于肉身，并将其置于心中。现在思想他在何种意义上宣称\BibleQuote{外面的人朽坏，内里的人却日日更新。}你当然不会主张他是指肉身从死亡那一刻开始，在其命定的永恒腐朽状态中经历的那种腐败；而是指它为了基督的名在生前直至死亡的过程中所经历的\Emph{磨损}，在烦扰的忧虑、磨难以及酷刑和迫害中。如今，内里的人当然必须藉圣神的推动而更新，藉着信德和圣洁日复一日地进步，\Emph{在此世}今生，而不是\Emph{在那里}复活之后，因为我们的更新不是日复一日的渐进过程，而是一次完成的圆满。你也可以从以下经文中得知这一点，那里\Emph{宗徒}说：\BibleQuote{我们这轻微短暂的苦难，为我们成就了极重无比、永远的荣耀；我们不是顾念所见的，}即我们的苦难，\BibleQuote{而是顾念那所不见的，}即我们的赏报：\BibleQuote{因为所见的是暂时的，所不见的是永远的。}\BibleRef{格后 4:17}因为外面的人所经历的苦难和伤害，他断言只配我们轻视，因为它们是轻微和暂时的；他更喜欢那些永恒的、也是看不见的赏报，以及那\Quote{荣耀的重量}，它将是肉身在此朽坏中所忍受劳苦的平衡。因此，这段经文的主题并非\Emph{他们}归给外面的人的那种腐败，即肉身的彻底毁灭，以否定复活。同样，他在别处也说：\BibleQuote{只要我们与\Emph{基督}一同受苦，也必与祂一同得荣耀；因为我算定现在的苦楚，不配与将来要显于我们的荣耀相比。}\BibleRef{罗 8:17}这里他又向我们表明，我们的苦楚小于它们的赏报。如今，既然我们是藉着肉身与\Emph{基督}一同受苦——因为肉身的特性就是因受苦而磨损——那么与基督一同受苦所应许的赏报也属于同一肉身。因此，当他将要指出苦难是肉身特别的责任时——根据他已经作出的声明——他说：\BibleQuote{我们到了\Place{马其顿}，我们的肉身没有安息；}\BibleRef{格后 7:5}然后，为使灵魂与身体同受苦难，他补充说：\BibleQuote{我们处处受困扰；外有争战，}这当然消磨肉身，\BibleQuote{内有恐惧，}这折磨灵魂。因此，尽管外面的人朽坏——不是指失去复活，而是指忍受磨难——从这\Emph{圣经}中我们可以理解，它并非没有内里的人而独自受苦。所以，两者将一同得荣耀，正如他们一同受苦一样。与他们一同受苦相伴的，也必然是他们一同分享赏报。\par

% structure: p0083 source=h2 target=section reason=relative-heading-rank; base=h2

\section{第四十一章 我们的帐棚解散与我们的身体复活一致}

他继续在同一经文中表达同样的情感，将报酬置于痛苦之上。他说：\BibleQuote{因为我们知道，如果我们地上帐幕式的房屋拆毁了，我们必得天主所造、非人手所成的永远居所。}\BibleRef{格后 5:1} 换言之，由于我们的肉身正因痛苦而解体，我们将在天上获得一个居所。他记起福音中（主所赐的）赏报：\BibleQuote{为义而受迫害的人是有福的，因为天国是他们的。}\BibleRef{玛 5:10} 然而，当他这样对比奖赏的报酬时，他并没有否认肉身的恢复；因为报酬应归于那承受解体的同一实体——当然就是肉身。然而，既然他曾称肉身为\Emph{房屋}，他愿意巧妙地用同一术语来比较最终的奖赏；应许那因痛苦而解体的房屋本身，通过复活得着一所更好的房屋。正如主也应许我们在祂父家有许多住处；\BibleRef{若 14:2} 尽管这或许可被理解为这世界的居所，在其结构解体时，天上应许了永恒的居所；然而，随后明显提及肉身的上下文似乎表明前面的言论并无此意。因为宗徒作出区分，继续说道：\BibleQuote{我们在这帐幕里叹息，渴望穿上我们那来自天上的住所；果然穿上，就不至于赤身了。}\BibleRef{格后 5:2} 意思是，在我们脱下肉身的衣服之前，我们愿穿上不朽的天上荣耀。这恩宠的特权等待着那些在主来临时仍活在肉身中的人，他们因着反基督时代的压迫，借着瞬间的死亡（通过突然的改变）而有资格加入复活的圣人；正如他写给得撒洛尼人的话：\BibleQuote{我们藉主的话告诉你们这件事：我们这些活着存留到主来临时的人，决不会在已死的人以前。因为主必亲自从天降来，有呼喊声、总领天使的声音、和天主的号声；那些在基督内死了的人先要复活；然后我们这些活着还存留的人，同时与他们一起被提到云中，到空中迎接主：这样我们就常与主同在。}\BibleRef{得前 4:15}\par

% structure: p0085 source=h2 target=section reason=relative-heading-rank; base=h2

\section{第四十二章 死亡改变而不毁灭我们可朽的身体；巨人的遗骸。}

他给格林多人解释这些人将要经历的转变，当他写道：\\BibleQuote{我们众人固然都要复活（但并非都要改变），在顷刻之间，一眨眼，在最后的号角响起时}——因为只有那些仍活在肉身中的人才会经历这种改变。\\BibleQuote{至于死者，}他说，\\BibleQuote{必要复活，\\Emph{并且}我们必要改变。}现在，经过仔细考量这既定的次序，你就能将后面的话与前面的文意对应起来。因为当他加上：\\BibleQuote{这必朽的必须穿上不朽的，这必死的必须穿上不死的，}\\BibleRef{格前 15:51}这必定就是那座我们切望穿上的、来自天上的居所，并在这现世的肉身中叹息——当然是指我们将最终被找到时的这肉身；因为他说我们在这帐棚中备受重负，我们确实不愿被剥去，反而宁愿被穿上\\Emph{在里面}，如此，必死性被生命吞灭，也就是说，当我们被转变时，穿上那来自天上的外衣。因为谁不愿意，当他仍在肉身中时，穿上不朽，并以快乐逃避死亡而延续生命？藉着必须代替死亡的转变，而不必遭遇那要追索最后一分钱的\\OriginalTerm{阴府}{Hades}？然而，那已经穿过\\Emph{阴府}的人，在复活后也注定要获得这种改变。因为由此我们就可以明确宣告，肉身\\Emph{绝对}要复活，并因即将到来的改变而拥有天使的形态。现在，如果只有那些在肉身中被找到的人才必须经历改变，以便必死性被生命吞灭——换言之，肉身被天上的永恒衣裳覆盖——那么，要么那些在死亡中被找到的人就得不到生命，因为他们被剥夺了生命的材料，即肉身；要么，这些人也必须经历改变，好使他们里面的必死性也被生命吞灭，因为他们也注定要获得生命。但你说，在死者身上，必死性已经被生命吞灭了。不，当然不是所有情况。因为有多少人很可能就是刚刚死去的人——如此新近被埋葬，以至于他们里面没有任何朽坏？因为你当然不会认为某物朽坏了，除非它被切断、消除、从我们的知觉中消失，全然不再显现。有古代巨人的尸骸；显然它们没有完全腐烂，因为它们的骨架仍然存在。我们已在别处谈及此事。例如，就在最近，就在本城，当有人亵渎地在许多古墓上奠基\\OriginalTerm{奥德翁}{Odeum}时，人们惊恐地发现，大约五百年后，骨头仍保持湿润，头发仍未失去香味。可以肯定，不仅骨头保持坚硬，而且牙齿也经久不腐——两者都是那要在复活中重新生长出来的身体的长久种子。最后，即使所有死者中的一切必死之物那时都已腐朽——无论如何被死亡、时间和年岁所吞噬——难道没有什么会被\\BibleQuote{被生命吞灭，}并且被不朽的衣裳覆盖和穿戴吗？现在，那说必死性将被生命吞灭的人，已经承认死者\\Emph{并没有被}那些其他\\Emph{前面提到的吞噬者}所毁灭。的确，一切都由天主的作为来成全和实现，\\Emph{并且}不是靠自然的法则，这将极为合适。因此，既然必死的必须被生命吞灭，它就必须显现出来以便被吞灭；也需要被吞灭，以便经历最终的转变。如果你说要点燃火焰，你就不能说有时需要、有时不需要点火之物。同样，当他插入这句话：\\BibleQuote{假如我们脱下之后，仍不被发现赤身，}\\BibleRef{格后 5:3}——当然是指那些在末日活着并在肉身中的主所找到的人——他并不是说，他所描述为赤身或脱去的人，在任何其他意义上算是赤裸，而是意味着他们应该被理解为重新穿上了他们所脱去的同一实体。因为虽然当他们肉身被搁置，或某种程度上撕裂或磨损时，他们会被发现赤裸（这种状况完全可以称为\\Emph{赤裸}），但他们以后会重新获得它，以便重新穿上肉身后，也能将不朽的超衣穿在上面；因为外衣除非穿在已经穿着的人身上，否则无法合身。\par

% structure: p0087 source=h2 target=section reason=relative-heading-rank; base=h2

\section{第四十三章 \\Person{圣保禄}称我们住在肉身时是远离主，这并不贬抑我们的教义。}

同样，当他说：\BibleQuote{因此我们时常放心大胆，并且深知，我们住在肉身内时，是远离主的；因为我们凭信德往来，并非凭目睹，}\BibleRef{格后 5:6} 显然，这句话并无诋毁肉身之意，仿佛肉身使我们与主分离。因为这里向我们明确提出一项劝诫，要我们轻视今世生命，因为我们只要在此世生活，就远离主——凭信德往来，而非凭目睹；换言之，凭望德，而非凭现实。因此他接着补充说：\BibleQuote{我们实在放心大胆，并认为更好是离开肉身，与主同住；}\BibleRef{格后 5:8} 目的就是使我们凭目睹而非凭信德往来，凭现实而非凭望德。请注意，他在这里也将蔑视肉身归于殉道的卓越。因为凡离开肉身的人，并不能立刻住在主的面前，除非凭借殉道的特权，他才能获得在乐园中的寓所，而非在下界。如今，宗徒是词穷而无法表达离世吗？还是他故意使用新奇的措辞？因为他想表达我们暂时离开肉身，就说我们是客旅，远离肉身，正如出外的人过些时候返回家园。然后他向众人说：\BibleQuote{因此我们无论身在何处，都切望讨天主喜悦，因为我们都必须出现在基督耶稣的审判台前。}\BibleRef{格后 5:9} 若我们都如此，那便是全人都如此；若是全人，那么我们的内在之人与外在之人——也就是说，我们的身体与灵魂，都无一例外。接着他继续说：\BibleQuote{为使每人，} 正如他所说，\BibleQuote{按他在肉身内所行的，或善或恶，领取相应的报应。}\BibleRef{格后 5:10} 现在我请问，你如何理解这段经文？你认为它的结构混乱、观念颠倒吗？问题在于，是身体将要承受的事，还是已经在身体内完成的事？若是身体所要承受的事，那么无疑暗示了身体的复活；若是已在身体内完成的事，也同样得出这个结论：因为报应当然必须由身体来支付，因为行为是由身体完成的。因此，宗徒从一开始的全部论证，便在这句结语中得到阐明，其中宣告了肉身的复活；并且应当以严格符合这结论的意义来理解。\par

% structure: p0089 source=h2 target=section reason=relative-heading-rank; base=h2

\section{第四十四章 圣保禄其他几处经文在一句话中的解释，以证实我们的教义。}

现在，你若查考提到外在之人和内在之人的经文之前的话语，岂不发现有关肉身的尊严与希望的全部真理？因为当他说：\BibleQuote{那吩咐光从黑暗中照耀的天主，曾照耀在我们心中，为使我们得知天主光荣的知识，这光荣反映在耶稣基督的面容上，}\BibleRef{格后 4:6} 并说\BibleQuote{我们把这宝贝放在瓦器中，}\BibleRef{格后 4:7} 显然指肉身，那么这肉身是要被毁坏，因为它只是\BibleQuote{瓦器}，出于尘土呢？还是要得荣耀，因为它是神圣宝贝的容器？如果那真光，即在基督内的光，本身含有生命，而这生命连同其光被托付给肉身，那么接受了生命的肉身怎能归于腐朽？当然，宝贝也会朽坏；因为可朽坏的东西托付给可朽坏的东西，如同把新酒装在旧皮囊里。当他接着说：\BibleQuote{我们身上时常带着耶稣\Emph{主}的死亡，}\BibleRef{格后 4:10} 那曾被称作天主的殿的实体，如今又可称作基督的坟墓吗？但我们为什么在身上带着主的死亡？为的是，如他所说，\BibleQuote{使\Emph{祂}的生命也得以彰显。} 在哪里？\BibleQuote{在身体上。} 在什么身体？\BibleQuote{在我们必死的身体上。} 因此是在肉身内，这肉身虽因罪而必死，却因恩宠而活着——当你看到目的是\BibleQuote{使基督的生命得以彰显在它内}时，这恩宠何其大！那么，在一个与救恩无份、不断消解的实体上，基督的生命（永恒、连续、不朽坏且已经是天主的生命）将要彰显吗？否则，那要在我们身体上彰显的主的生命属于哪个时期？那一定是祂在受难前所活的生命，这生命不仅曾公开显于犹太人中间，如今也向万民显明了。因此，那生命是指\Quote{打破了死亡的砥柱和阴间的铜闩}的生命——这生命从此以后是我们的，也将永远是我们的。最后，它要在身体上彰显。何时？死后。如何？藉着在我们身体内复活，正如基督也\Emph{在祂的身体内}复活。但恐怕有人会反驳说，耶稣的生命甚至现在也要在我们身上通过圣德、忍耐、正义和智慧的生活彰显，这些在主的生活中非常丰富；宗徒最富远见的智慧便加入了这一目的：\BibleQuote{因为我们活着的人，时常为耶稣的缘故被交于死亡，为使耶稣的生命也彰显在我们必死的身体上。}\BibleRef{格后 4:11} 因此，即使当我们死了，他说这事要发生在我们身上。若是如此，这除了在我们复活后的身体上还能如何成就？因此他在结语中补充说：\BibleQuote{我们知道那使\Emph{主}耶稣复活的那位，也必要使我们与\Emph{祂}一同复活。}\BibleRef{格后 4:14} 祂已从死者中复活了。但或许\BibleQuote{与\Emph{祂}一同}意思是\BibleQuote{像\Emph{祂}一样}：那么，若是像\Emph{祂}一样，当然就不能没有肉身。\par

% structure: p0091 source=h2 target=section reason=relative-heading-rank; base=h2

\section{第四十五章 圣保禄所说的旧人与新人解释。}

但他们在盲目中，再次在旧人与新人的问题上自陷困境。当宗徒命令我们\BibleQuote{脱去那个依照欺骗情欲而败腐的旧人，并在我们心思的灵里得以更新，并且穿上那个按天主在正义和真正的圣德中所创造的新人}\BibleRef{弗 4:22}时，（他们主张）这样在两种实体之间作出区分，\Emph{并将旧人归于肉体，新人归于精神}，他就把一种永久的败腐归于旧人——即肉体。如今，如果你按照实体的先后顺序，灵魂不可能是新人，因为它在两者中是后出现的；肉体也不可能是旧人，因为它是先存在的。因为天主的\Emph{创造}的手和祂的\Emph{嘘气}之间，间隔了多短的时间呢？我敢说，即使灵魂远在肉体之前，正是由于灵魂必须等待使自己得以完全，这就使得另一个实际上成了先者。因为凡给一项工作带来最终完成和完美的事物，尽管在纯粹次序上它是在后的，但在效果上却有优先性。更何况那没有它，先前的诸物便不可能存在的事物，更是优先的。如果肉体是旧人，它是在何时成为这样的？从一开始吗？但亚当是完全的新人，而那个新人不可能有任何部分是旧人。而从那时起，自从对人繁衍的祝福被宣告以来\BibleRef{创 1:28}，肉体和灵魂就有了同时的诞生，时间上没有任何可计算的差异；以至于两者甚至一起在母腹中被生成，正如我们在\Emph{\WorkTitle{论灵魂}}中所指出的。在母腹中同时，它们在出生上时间上也相同。两者无疑是由属于两种实体的人类父母所生，但不是在两个不同时期；相反，它们如此完全合一，以至于\Emph{在时间上}无一方先于另一方。更正确的说法是，我们要么完全是旧人，要么完全是新人，因为我们无法理解如何可能成为别的。但宗徒提到了旧人的一个非常清楚的标志。因为他说：\BibleQuote{脱去……照着从前的生活方式，旧人}\BibleRef{弗 4:22}，他并不是\Emph{说}关乎任何一种实体的年长。他命令我们脱去的确实不是肉体，而是他在另一处经文中所指明为\BibleQuote{肉性的作为}的事\BibleRef{迦 5:19}。他对人的身体并不加以谴责，反而写下了这样的话：\BibleQuote{你们既然脱去了谎言，就要各人与邻舍说实话，因为我们彼此是肢体的肢体。你们生气，却不要犯罪：不可让太阳在你们的愤怒中落下，也不可给魔鬼留地步。偷窃的，不要再偷；反倒要劳力，亲手作正经事，使他可以分给那有需要的人。败坏的话一句都不可出你们的口，只要说那有利于信德建设、使听者得恩宠的话。不要叫天主的圣神担忧，你们原是受了他的印记，等候得赎的日子。一切苦毒、恼怒、忿怒、喧嚷、毁谤，连同一切的恶毒，都当从你们中间除掉；并要以恩慈相待，存怜悯的心，彼此饶恕，正如天主在基督里饶恕了你们一样。}\BibleRef{弗 4:25} 既然如此，那些以为肉体是旧人的人，为什么不赶紧自己去死，好藉着脱去旧人来满足宗徒的训诫呢？至于我们自己，我们相信整个信德要在\Emph{肉体中}实行，甚至要通过\Emph{肉体}来实行，因为肉体有口可以宣讲一切圣善的言语，有舌可以避免亵渎，有心可以避免一切恼怒，有手可以劳作和施舍；同时我们也主张，旧人和新人都与道德行为的差异有关，而与本性的任何不同无关。正如我们承认，那按从前的生活方式是\BibleQuote{旧人}的，也是败腐的，并按照\BibleQuote{欺骗的情欲}得了它的名字，同样（我们主张）那\BibleQuote{旧人是就从前的生活方式而言}\BibleRef{弗 4:22}，而非就肉体而言，因为肉体并无永久的解体。而且，它（指旧人旧性）在肉体中仍未受损，并且在那本性上相同，甚至当它已成为\Quote{新人}时也是如此；因为它所脱去的是它的罪恶生活历程，而不是它的物质实体。\par

% structure: p0093 source=h2 target=section reason=relative-heading-rank; base=h2

\section{第四十六章 \Person{圣保禄}常谴责的是\OriginalTerm{肉性的作为}{works of the flesh}，而非肉体的实体。}

你们或许注意到，宗徒处处如此谴责肉身的行为，以至于似乎是在谴责肉身本身；但无人能以为他有这样的看法，因为他接着提出了另一种含义，尽管有几分相似。因为当他明确宣称\Quote{凡属于肉身的，不能中悦天主}时，他立刻将这话从异端的意义挽回至健全的意义，补充说：\Quote{但你们不属于肉身，而是属于圣神。}\BibleRef{罗 8:8}如今，他否认那些显然属于肉身的人属于肉身，表明他们并非活在肉身的行为中，因此不能中悦天主的，不是那些属于肉身的，而是那些随从肉身生活的人；相反，那些虽然处于肉身之中，却随从圣神行事的人，却中悦天主。他又说：\Quote{身体是死的}；但这是\Quote{因为罪}，而\Quote{圣神却是生命，因为正义}。然而，当他如此将生命与肉身中的死亡对立时，他无疑将正义的生命应许给了那同一状态——他为之确定了罪之死亡的状态。但若生命不在他与之对立的那事物所在之处——即那要从肉身中除去显然的死亡——那么他所做的\Quote{生命}与\Quote{死亡}的对立就毫无意义。如果生命如此从肉身中除去死亡，它只有渗透到它所排除的事物所在之处，才能做到这一点。但我何必诉诸棘手的论证呢？宗徒本身已极其明白地处理了这主题。因为他说：\Quote{如果那使耶稣从死者中复活者的圣神住在你们内，那么那使耶稣从死者中复活的，也必藉着住在你们内的圣神，使你们有死的身体活过来}\BibleRef{罗 8:11}；因此即使有人假定灵魂是\Quote{有死的身体}，他也（既然他无法否认肉身也同样如此）不得不承认肉身因其参与同一状态而得以恢复。此外，从以下话语你可以学到，被谴责的是肉身的行为，而非肉身本身：\Quote{因此，弟兄们，我们并不是欠\Emph{肉身}的债，去随从肉身生活；如果你们随从肉身生活，你们必要死；但你们若藉着圣神，治死身体的作为，必要生活。}现在（让我分别回答每一点），既然救恩应许给那些生活在肉身中、却随从圣神行事的人，那么与救恩为敌的就不是肉身，而是肉身的行为。然而，当这肉身的行为——它是死亡的原因——被除去时，肉身就显为安全，因为它已脱离了死亡的原因。因为他说：\Quote{在基督耶稣内生命之圣神的法律，已使我脱离罪与死亡的法律}——那正是他先前提到住在我们肢体中的法律。因此，我们的肢体不再受死亡的法律管辖，因为它们不再服事罪的法律，从\Emph{两者}中它们已获得自由。\Quote{因为法律因肉身的软弱所不能行的，天主却做到了：祂派遣了自己的儿子，带着罪恶肉身的形状，并因着罪，在肉身上定了罪}\BibleRef{罗 8:3}——而不是在罪中定了肉身的罪，因为房屋不应与其住户同被定罪。他确实说过\Quote{罪住在我们身体内}\BibleRef{罗 7:20}。但罪的定罪是肉身的开释，正如其不定罪使肉身服从罪与死亡的法律。同样，他称\Quote{肉性思念}先是\Quote{死亡}\BibleRef{罗 8:6}，之后又说是\Quote{对天主的仇视}；但他从未将这些归于肉身本身。但你会说：那么肉性思念应归于什么呢，若不是归于肉性的实体本身？我同意你的反对，只要你证明给我看，肉身有它自己的辨别力。然而，若肉身没有灵魂便什么也不能认识，你必须明白，肉性思念应归于灵魂，尽管有时归之于肉身，因为它是为肉身并通过肉身而被服务的。因此（宗徒）说\Quote{罪住在肉身中}，因为挑起罪的灵魂暂时居住在肉身中，这肉身确实注定死亡，但不是因它自己，而是因罪。因为他在另一处也说：\Quote{你们怎么仍像生活在世俗中一样行事呢？}\BibleRef{哥 2:20}这里他并非写给死者，而是写给那些本应停止按世俗方式生活的人。\par

% structure: p0095 source=h2 target=section reason=relative-heading-rank; base=h2

\section{第47章 圣保禄始终应许肉身永生}

因为那旧人（正如他所说）是\BibleQuote{与基督一同被钉在十字架上} \BibleRef{罗 6:6}，它必须在世界之后活着，这不是指肉身结构，而是指道德行为。此外，如果我们不这样理解，那么（无论如何）被钉穿的不是我们的身体框架，我们的肉身也没有承受基督的十字架；而是他所附加的意义：\BibleQuote{为使罪身失效}，通过生活的改正，而非本质的毁灭，正如他继续说：\BibleQuote{使我们从今以后不再服事罪恶}；并且我们应当相信自己\BibleQuote{与基督同死}，以至于\BibleQuote{我们也必与祂同活}。按照同一原则，他说：\BibleQuote{同样，你们也当看自己是死的。}向着什么？向着肉身？不，而是向着\BibleQuote{罪}。因此，就肉身而言，他们将被拯救——\BibleQuote{在基督耶稣内，活于天主}，当然是通过肉身，他们并不向着肉身死；因为他们向着\BibleQuote{罪}死，而不是向着肉身。因为他进一步阐述说：\BibleQuote{所以，不要让罪在你们可朽的身体内为王，使你们顺从它，并将你们的肢体献给罪，作不义的器具；却要将你们自己献给天主，如同从死中复活的人}——不仅仅是活着，而是如同从死中复活——\BibleQuote{并将你们的肢体献给义，作义的器具。}又说：\BibleQuote{你们怎样将肢体献给不洁和不法作奴隶，以致于不法，现今也要照样将肢体献给义作奴隶，以致于圣德。因为你们作罪恶的奴隶时，不受正义的约束。那时你们有什么果实呢？那些事如今你们引以为耻，因为它们的结局就是死亡。但现今你们脱离了罪恶，作了天主的奴隶，就有果实以致于圣德，结局就是永生。因为罪恶的薪俸是死亡，但天主的恩赐是在我们的主基督耶稣内的永生。}因此，贯穿这一系列经文，他既将我们的肢体从不义和罪恶中撤回，又将它们应用于正义和圣德，并将它们从死亡的薪俸转移至永生的恩赐，便无疑将救恩的赏报应许给肉身。然而，如果这种修行的报酬不是肉身所能企及的，那么特别为肉身规定任何圣洁和正义的准则就完全不合情理；同样，如果圣洗的重生不是开创一条趋向肉身恢复的道路，那么圣洗也不可能被恰当地为肉身而制定。宗徒本人暗示这一思想说：\BibleQuote{难道你们不知道，我们这受过洗归于基督耶稣的人，是受洗归于祂的死吗？所以我们借着洗礼与祂同葬于死，为的是基督怎样从死者中复活，我们也怎样在新生活中行走。}\BibleRef{罗 6:3}为使你们不以为这只是指我们借着圣洗、因着信德所应行走的新生活，宗徒以极高的预见补充说：\BibleQuote{我们若在祂死的形状上与祂联合，也要在祂复活的形状上与祂联合。}我们在圣洗中借比喻死去，但我们在肉身中真实地复活，正如基督一样，\BibleQuote{好使罪恶怎样在死亡中为王，恩宠也怎样借着正义而为王，使人借着我们的主耶稣基督获得永生。}\BibleRef{罗 5:21}但若不也在肉身中，怎能如此？因为死亡在哪里，死亡后的生命也必在哪里，因为生命原先就在那里，死亡后来才在那里。如今，若死亡的权势仅作用于肉身的解体，同样，死亡的对立面——生命，也应当产生相反的效果，即肉身的恢复；如此，正如死亡以其势力吞噬了肉身，这必朽的被不朽吞噬之后，肉身也能听到向死亡发出的挑战：\BibleQuote{死亡啊，你的刺在哪里？坟墓啊，你的胜利在哪里？}\BibleRef{格前 15:55}因为如此，\BibleQuote{罪恶在哪里增多，恩宠在那里也格外丰富。}\BibleRef{罗 5:20}同样，\BibleQuote{有力量在软弱中得以成全。}\BibleRef{格后 12:9}——拯救失落的，复活死亡的，医治受伤的，治愈衰弱的，赎回迷失的，释放被奴役的，唤回迷路的，扶起跌倒的；这一切从地上到天上，正如宗徒教训斐理伯人所说的：\BibleQuote{我们却是天上的国民，并且等候主耶稣基督从天上降临，祂要改变我们卑贱的身体，使它相似祂荣耀的身体}\BibleRef{斐 3:20}——当然是在复活之后，因为基督自己也没有在受苦之前受光荣。这些必定是他恳求罗马人\BibleQuote{献上}为\BibleQuote{活的祭品，圣洁的，中悦天主}的\BibleQuote{身体}。\BibleRef{罗 12:1}但若这些身体将要朽坏，怎能是\Emph{活的}祭品？若它们被亵渎地玷污，怎能是\Emph{圣洁的}？若它们被定罪，怎能是\Emph{中悦天主的}？来吧，如今告诉我，那些躲避圣经之光的异端者是如何理解致得撒洛尼人书中的那段经文——因其清晰，我甚至以为它是用阳光写成的：\BibleQuote{愿和平的天主亲自全然圣化你们。}而且，仿佛这还不够明白，它接着说：\BibleQuote{并愿你们的整个身体、灵魂和心神，在我们的主耶稣基督来临时，保持无可指摘。}\BibleRef{得前 5:23}在此，你看到人的整个本质注定得救，而且\Emph{那}是在主来临时，别无其他时刻，这正是复活的关键。\par

% structure: p0097 source=h2 target=section reason=relative-heading-rank; base=h2

\section{第四十八章 解释有关死人复活的伟大章节中的若干经文，以捍卫我们的教义。}

但是，你们说，\BibleQuote{血肉}，你们说，\BibleQuote{不能继承天主的国。}\BibleRef{格前 15:50}我们很知道这也是经上所载；但是，虽然我们的对手把它摆在战斗的前线，我们故意保留这个异议直到现在，以便我们在最后一次进攻中推翻它，在我们清除了所有辅助它的疑问之后。然而，他们必须设法现在甚至回想我们前面的\Emph{论证}，以便最初暗示这段经文的情景有助于我们判断得出其含义。依我看来，宗徒为格林多人陈明了他们教会纪律的细节，总结了他自己福音的实质，以及他们在阐明主的死亡和复活时的信仰，目的是为了从中推论出我们望德的规则及其基础。因此他附加了这句话：\BibleQuote{如果基督被传讲已经从死人中复活，你们中间有些人怎么说没有死人的复活呢？如果没有死人的复活，基督也就没有复活；如果基督没有复活，我们的宣讲就是徒然，你们的信德也是徒然。并且，我们被发现在天主面前作假见证，因为我们曾为天主作证说祂使基督复活了，而祂并没有使祂复活，倘若死人不会复活的话。因为如果死人不会复活，基督也就没有复活；如果基督没有复活，你们的信德就是徒然，因为你们仍在你们的罪中，那些在基督内睡了的人也灭亡了。}\BibleRef{格前 15:12}那么，他在整段经文中明显努力让我们相信的是什么？死人的复活，您说，那是被否认的：他当然希望我们因他所举的例子——主的复活——而相信它。当然，您说。那么，例子是从不同的情况借用来的，还是从类似的情况？从类似的情况，毫无疑问，是您的回答。那么基督是如何复活的？在肉身内，还是不在？毫无疑问，既然你们被告知祂\BibleQuote{按照圣经死了}，并且\BibleQuote{祂被埋葬，\Emph{按照圣经}}，无非是在肉身内，你们也会承认祂是在肉身内从死人中复活的。因为同一个在死亡中倒下的、躺在坟墓里的身体也复活了；（这与其说是基督在肉身内，不如说是肉身在基督内。）因此，如果我们必须按照基督的榜样复活，祂在肉身内复活了，我们肯定不会按照那个榜样复活，除非我们自己也在肉身内复活。\BibleQuote{因为，}他说，\BibleQuote{既然死亡是由人而来的，死人的复活也是由人而来的。}（他这样说）一方面是为了区分两个作者——死亡来自亚当，复活来自基督；另一方面是为了使复活作用于与死亡相同的实体上，通过以\Emph{人}这个称号来比较作者本身。因为如果\BibleQuote{在亚当内众人都死了，照样在基督内众人也都要复活，}\BibleRef{格前 15:22}他们在基督内的复活必须是在肉身内，因为他们在亚当内的死亡是在肉身内发生的。\BibleQuote{但各人要按着自己的次序，}因为当然\Emph{各人也将是}各人\Emph{在自己的身体内}。因为\Emph{次序}将分别安排，鉴于个人的功绩。既然功绩必须归于身体，那么次序也必须按照身体来安排，以便与它们的功绩相关。但是，既然\BibleQuote{有些人也为死人受洗，}我们将看看是否有好的理由这样做。现在，可以肯定的是，他们采用这种（做法）是基于这样的假设，即认为（有疑问的）代洗将有益于另一个人的肉身，以期复活；因为除非是身体的\Emph{复}\Emph{活}，否则这种肉体洗礼的程序不会保证什么。\BibleQuote{那么他们为什么为死人受洗呢？}他问道，除非这些受洗的身体复活？因为不是灵魂被洗礼的水圣化：它的圣化来自\BibleQuote{回应}。\BibleRef{伯前 3:21}\BibleQuote{为什么，}他问，\BibleQuote{我们每时刻都处于危险中？}\BibleRef{格前 15:30}——当然是指通过肉身。\BibleQuote{我天天冒死，}（他说）；那无疑是在身体的危险中，在那危险中\BibleQuote{他甚至曾在厄弗所与野兽战斗，}——甚至与那些在亚细亚带给他如此危险和麻烦的野兽战斗，他在写给同一教会\Emph{格林多}的第二封书信中暗示了这一点：\BibleQuote{弟兄们，我们不愿你们不知道我们在亚细亚所遭遇的患难，我们被压太重，力不能胜，甚至连活命的指望都绝了。}\BibleRef{格后 1:8}现在，如果我没有弄错，他列举所有这些细节，是为了他不愿意让他在肉身中的冲突被认为是无用的，他可能引出对肉身复活的不动摇的相信。因为那冲突（在身体内承受的）必须被视为无用，如果它没有复活的前景。\BibleQuote{但有人会说，死人怎样复活？以什么样的身体来呢？}\BibleRef{格前 15:35}现在他讨论身体的品质，人们要重新取回的究竟是同一身体还是不同的身体。然而，既然这样一个问题必须被视为后续问题，我们顺便指出，复活被确定为身体的复活，甚至从这一点就可以看出：问题正是关于\Emph{身体}的品质的。\par

% structure: p0099 source=h2 target=section reason=relative-heading-rank; base=h2

\section{第四十九章 同一主题继续。宗徒从死人中排除什么？确实不是肉身的实体。}

我们现在来到整个问题的核心：宗徒究竟将哪些实体、以及何种性质的实体排除于天主的国之外？前面的论述也给了我们关于这一点的线索。他说：\BibleQuote{第一个人是出于地，属于土}——即由尘土所造，就是亚当；\BibleQuote{第二个人是出于天} \BibleRef{格前 15:47}——即天主的圣言，就是基督，然而祂之所以是\Emph{人}（虽说是\Quote{出于天}），无非是因为祂本身是血肉和灵魂，正如任何一个人一样，正如亚当一样。事实上，在先前的一段经文中，祂被称为\BibleQuote{第二个亚当}，这同一名称来自祂对实体之分有，因为就连亚当也不是出于人的种子而成血肉，而基督在这方面也与祂相似。\BibleQuote{那属于土的怎样，凡属于土的也怎样；那属于天的怎样，凡属于天的也怎样。} \BibleRef{格前 15:48} \Emph{这样}（他是说）指实体？还是首先指训练，然后指那训练所旨在获得的尊荣与价值？然而，决不是指实体，属土的和属天的将会被分开，因为它们已被宗徒一劳永逸地指称为\Emph{人}。因为即使基督是唯一真正的\Quote{属天的}，甚至属于超天界的存有，祂仍然是具有身体和灵魂的人；一点也没有因祂分有双重实体的状况而与\Quote{属土}的特质分离。同样，那些在祂之后成为属天的人，也被理解为并非因他们现在的本性，而是因他们将来的光荣而被预指为具有这属天的特质；因为在前面的句子中，这关于尊荣差异的区分起先被引出时，曾显示出\BibleQuote{在天上的形体上有一种光荣，在地上的形体上有另一种光荣} \BibleRef{格前 15:40} ——\BibleQuote{太阳有一种光荣，月亮有另一种光荣，星辰也有另一种光荣；因为星辰与星辰在光荣上也有差异}，尽管在实体上没有差异。然后，在这样预先提出那现在就要追求、最后将享受的价值或尊荣的差异之后，宗徒加上了一个劝勉，要我们既在此生训练中效法基督的榜样，在那里达到祂光荣的高度：\Quote{我们怎样带了那属于土的肖像，也要怎样带那属于天的肖像。}我们确实带了那属于土的肖像，藉着分有他的过犯，藉着分有他的死亡，藉着被逐出乐园。如今，虽然亚当的肖像是在肉身上被我们带着，但我们并不是被劝告脱去肉体；若不是肉体，那就是行为方式，为使我们能将自己身上带那属于天的肖像——不再是\Emph{天主的肖像}，也不再是那在天上的存有的肖像；而是照基督的样式，藉着我们在此生走在圣洁、正义和真理中。他在整个这一段落中如此专注于道德行为的教导，以至于他告诉我们应当在这肉身上、在这训练与纪律的时期中带基督的肖像。因为当他说\Quote{\Emph{让我们带}}（使用祈使语气）时，他的话适用于现世生活，在此生活中人只以血肉和灵魂的实体存在；或者即使那是另一种实体，即属天的实体（我们的信德所期待的），那应许也是赐给那奉命要勤勉努力以配得奖赏的\Emph{实体}的。既然他使属土的和属天的肖像都在于道德行为——一个要弃绝，一个要追求——然后便一致地加上：\BibleQuote{因为我说这话}（即因我前面所说的，因为连词\BibleQuote{\Emph{因为}}连接了后面与前面的话）\BibleQuote{血肉不能继承天主的国} \BibleRef{格前 15:50}——他意思是血肉只能按前面提到的\Quote{属于土的肖像}来理解；既然这被认为在于\Quote{旧生活}，而这旧生活不接受天主的国，因此血肉因不接受天主的国而被归结为旧生活之\Emph{生命}。当然，既然宗徒从未将实体与人的\Emph{行为}混同，他在这里就不能使用这样的结构。然而，既他已论及那些仍在肉身中活着的人，说他们\BibleQuote{不在肉身中} \BibleRef{罗 8:9}，意思是不活在血肉的行为中，你就不该颠覆其形式或其实体，而只该颠覆那在实体（肉体）中所行、使我们与天主的国隔绝的行为。在向迦拉达人显示这些有害的行为之后，他声称要预先警告他们，正如他\BibleQuote{从前告诉过他们，行这样事的人不得继承天主的国} \BibleRef{迦 5:21}，甚至因为他们没有带属天的肖像，正如他们带了属土的肖像一样；因此，由于他们的旧生活，他们只能被视为无非是血肉。但即使宗徒突然抛出那句血肉必须从天的国被排除的句子，而没有任何预先的提示，难道我们不也同样有责任将这两种实体理解为那完全耽于血肉的老旧的人——即耽于吃喝，其中一项特征就是反对复活的信德：\BibleQuote{让我们吃喝吧，因为明天我们就要死了。} \BibleRef{格前 15:32} 当宗徒以插入方式加上这一句时，他谴责了血肉因其在吃喝中的享乐。\par

% structure: p0101 source=h2 target=section reason=relative-heading-rank; base=h2

\section{第五十章 血肉在何种意义上被排除于天主的国之外。}

然而，撇开所有这些指责血肉之作为的诠释，请允许我要求复活正是这些实体，且仅按其自然意义来理解。因为直接被否定给血肉的并非复活，而是天主的国——这国是复活之后的事（因为还有审判的复活）；而且，每一次特别复活被排除时，反而是对一般肉体复活的确认。既然清楚说明了复活不会导向何种状态，那么它就导向何种状态也就明白了；因此，既然按\Emph{人的}功过，复活中的区别是由他们在肉体内的行为造成的，而非由肉体的实体本身，那么即使从这一点也可以看出，血肉被排除在天国之外是因为它们的罪，而非因为它们的实体；而且，虽然就它们的自然状态而言，它们必为审判而复活，因为它们没有为天国而复活。再者，我要说：\Quote{血肉不能继承天主的国；}\BibleRef{格前 15:50}并且，宗徒正是公义地宣告了这一点：它们独自、就其本身而论，是为要表明还需要圣神（使它们有资格）进入天国。因为\Quote{那使人活的是圣神}为\Emph{我们}进入天主的国；\Quote{肉一无所益。}\BibleRef{若 6:63}然而，另有一些东西能为之获益，那就是圣神；并且借着圣神，还有圣神的工程。因此，血肉必须在任何情况下都同样按其固有的性质复活。但是那些蒙恩进入天国之人的，必须穿上不朽坏、不死的生命的能力；因为没有这个，\Emph{或}在能够获得它之前，他们不能进入天主的国。因此，正如我们已经说过的，血肉自身无法获得天主的国，这是有充分理由的。然而，既然\Quote{这可朽坏的（即肉）必须穿上不朽坏，这可死的（即血）必须穿上不死}\BibleRef{格前 15:53}，借着复活后将发生的变化，因此有最充分的理由：血肉在那变化和穿着之后，将能够继承天主的国——但不可没有复活。有人主张，短语\Quote{血肉}由于割损礼而意指犹太教，后者本身与天主的国相距甚远，被视为\Quote{旧日或先前的行事}，并且在宗徒的另一段经文中也被以此名称呼，他写信给迦拉达人时说：\Quote{当天主乐意将祂的儿子启示给我，叫我在外邦人中传扬祂，我就没有同血肉商量}（此话意指）割损礼，即犹太教。\par

% structure: p0103 source=h2 target=section reason=relative-heading-rank; base=h2

\section{第51章 \Person{耶稣}以其成肉身之性坐在天主右边，乃是我们肉身复活的保证。}

那个我们留作结论的论点，现在将成为为所有人辩护的理由，也为宗徒本人辩护；宗徒本人如果像有些人所说、那样草率地、盲目地、不加区分地、无条件地将一切血肉排除在天主之国之外，甚至排除在天堂宫廷之外，那么他确实会犯下极大的轻率之罪。因为耶稣仍坐在圣父的右边，\BibleRef{谷 16:19} 是人，却是天主——末后的亚当，\BibleRef{格前 15:45} 又是原始的圣言——是血肉，却比我们的更为纯洁——祂\Quote{要怎样升\Emph{到天上}，也要怎样降来} \BibleRef{宗 1:9}，本质和形貌都相同，正如天使所证实的，甚至被刺透祂的人所认出。祂被称为\Quote{天主与人之间的中保}，亲自保管着由双方交托给祂的血肉存款——就是整个完美的抵押和保证。因为正如\Quote{祂赐给我们圣神的抵押} \BibleRef{格后 5:5}，祂也接受了我们血肉的抵押，并将它带到天上，作为将来要完全恢复的担保。血肉啊，不要为任何忧虑而烦恼；在基督内，你们已经获得了天堂和天主之国。否则，如果有人说你们不在基督内，就让他们也说基督不在天上，因为他们已经否认了你们的天堂。同样，他说：\Quote{朽坏也不能承受不朽坏的。} \BibleRef{格前 15:50} \Emph{他说这话}，并不是要你们把血肉当作朽坏，因为它们本身反而是朽坏的主体——我指的是通过死亡，因为死亡与其说是腐化我们的血肉，不如说是消耗它们。但既然他已经明确说过，血肉的行为不能承受天主之国，为了加强这一说法，他剥夺了朽坏本身——即死亡，它从血肉的行为中受益匪浅——对不朽坏的任何继承。因为稍后，他描述了死亡本身的死亡：他说：\Quote{死亡被胜利吞灭。死亡啊，你的刺在哪里？坟墓啊，你的胜利在哪里？死亡的刺就是罪}——这里就是\Emph{朽坏}；\Quote{而罪的力量就是法律} \BibleRef{格前 15:54}——毫无疑问，就是那另一条法律，他曾描述为\Quote{在他肢体中交战，反抗他心思的法律} \BibleRef{罗 7:23}——当然指的是违背他意愿的犯罪能力。他在前面的段落（格林多书）中说，\Quote{最后被消灭的仇敌是死亡。} \BibleRef{格前 15:26} 这样，朽坏就不能承受不朽坏；换句话说，死亡不再存留。它何时并如何停止呢？在\Quote{一刹那，一眨眼，在最后的号角吹响时，死人要复活成为不朽坏的。} 这些人不就是先前朽坏的人——即我们的身体，也就是我们的血肉吗？我们要改变。但改变成什么状态，如果不是我们将来被找到的状态呢？\Quote{因为这必朽坏的必须穿上不朽坏，这必死的必须穿上不死。} 这必死的是什么，如果不是血肉？这必朽坏的是什么，如果不是血？此外，为了使你不以为宗徒有其他意思，在他谆谆教诲你时，并为了使你明白他认真地将他的话应用于血肉，当他说\Quote{\Emph{这}必朽坏的}和\Quote{\Emph{这}必死的}时，他是在触摸自己的身体表面时说出这些话的。他当然不可能说出这些短语，除非指一个可触摸、可见的对象。这种表达暗示了身体的展示。此外，朽坏的身体是一回事，朽坏是另一回事；同样，必死的身体是一回事，必死是另一回事。因为受苦的是一回事，使之受苦的是另一回事。因此，那些受朽坏和必死支配的东西，即血肉，也必然能够接受不朽坏和不死。\par

% structure: p0105 source=h2 target=section reason=relative-heading-rank; base=h2

\section{第五十二章 从圣保禄的种子类比，我们得知已死的身体将复活，并饰以永生之器具。}

现在让我们看看他断言死去的人将以什么身体来临。他巧妙地步步进逼，立即阐明这一点，仿佛有一个诘难者用某个问题刁难他。\Quote{糊涂人哪，}他说，\Quote{你所种的，若不先死，决不会活起来。}\BibleRef{格前 15:36}从这个\Emph{种子}的例子显然可见，只有那经历过死亡的肉身才被赋予生命，因此其余的一切问题都变得足够清楚。因为任何与例子所提示的想法不相容的事都不能被理解；也不允许你们从随后的句子\Quote{你所种的，不是那将要生出的身体}中推测，在复活中，将有一个不同的身体从那在死亡中种下的身体兴起。否则，你们就脱离了例子。因为如果麦子被种在地里并分解，大麦不会生长。不过，它在种类上并非完全相同；它的性质、品质和形态也不相同。那么，如果不是完全相同，它从\Emph{何来}？因为连腐烂本身也是那东西本身\Emph{的证明}，因为它是实际\Emph{谷粒}的\Emph{腐烂}。然而，宗徒本人不也提示了\Quote{那将要生出的身体}为何不是所种下的身体吗？他岂不是说：\Quote{不过是赤裸的颗料，可能是麦子，也可能是别的什么谷粒；但天主随自己的心意给它一个身体。}\Emph{给予}它——当然是指那被他称为赤裸地种下的谷粒。毫无疑问，你会说。那么，谷粒是足够安全的，天主将给它一个身体。但如果它不存在于任何地方，如果它不再复活，如果它真实自身不再复活，它如何安全？如果它不再复活，它就不安全；\Emph{而且}，如果它甚至不安全，它就不能从天主那里接受身体。但有一切可能的证据证明它是安全的。因此，天主为什么将给它\Quote{一个身体，随自己的心意}，即使它已经有了自己的\Quote{赤裸}的身体，除非是在它的复活中它不再赤裸？因此，那将是附加的物质，覆盖在\Emph{赤裸}的身体之上；而那被上层物质放置其上的东西绝没有消灭——反而增加了。然而，那接受增长的东西是安全的。事实是，它被种下时是最赤裸的谷粒，没有外壳覆盖，甚至没有发芽的穗，没有芒顶的保护，没有茎秆的荣光。然而，它从垄沟里长出，因丰盛的收成而充实，以紧凑的结构建造，以\Emph{美丽}的顺序构成，因栽培而坚固，四面披覆。这些情况使它成为来自天主的另一个身体，它变化不是由于消灭，而是由于扩大。天主给每一种子\Emph{赋与它自己的身体}——\BibleRef{格前 15:38}——当然不是指它初始身体意义上的\Quote{自己的}；以便它从外部从天主得来的东西最终也被算作它自己的。因此，紧紧抓住这个例子，并将其作为肉身之事的镜子牢记在心：相信那曾经在\Emph{死亡}中种下的同一\Emph{肉身}，将在\Emph{复活的生命}中结出果实——本质相同，只是更丰满完美；不是另一个，尽管以另一种形式出现。因为它将在自身内接受天主乐意按其功劳赋予它的恩宠和装饰。毫无疑问，正是在这个意义上，他说：\Quote{凡肉身不都是一样的}——意思不是否认本质的共性，而是否认特权的平等——将身体降低为荣誉的差异，而非本性的差异。本着这一观点，他以比喻的方式添加了一些动物和天体的例子：\Quote{人的肉身是一样}（即，天主的仆人，但确实是人），\Quote{走兽的肉身又是一样}（即，外邦人，先知实际上论他们说：\Quote{人如同无知的畜牲}），\Quote{鸟的肉身又是一样}（即，那些试图飞向天堂的殉道者），\Quote{鱼的肉身又是一样}（即，那些被圣洗之水淹没的人）。\BibleRef{格前 15:39}同样地，他以天体为例：\Quote{太阳的光荣是一样}（即，基督的光荣），\Quote{月亮的光荣另是一样}（即，教会的光荣），\Quote{星辰的光荣另是一样}（即，亚巴郎子孙的光荣）。\Quote{因为星与星的光荣不同：同样，也有地上的身体和天上的身体}（即，犹太人和基督徒）。\BibleRef{格前 15:41}现在，如果这些话不按比喻理解，那他对比骡子和鸢的肉身，以及天体和人的身体，就够荒谬了；因为它们在状况方面无法比较，在获得复活方面也无法比较。然后，最后，通过他的例子确凿地表明差异在于光荣而非本质，他补充说：\Quote{死人的复活也是这样。}怎么回事？无非是在光荣上有所不同。因为，再次论及同一本质的复活，并回到（他的）谷粒（比喻），他说：\Quote{播种的是可朽坏的，复活起来的是不可朽坏的；播种的是羞辱的，复活起来的是光荣的；播种的是软弱的，复活起来的是强健的；播种的是属生灵的身体，复活起来的是属神的身体。}当然，除了所种下的之外，没有别的东西复活；除了在土里腐烂的东西之外，没有别的东西被种下；而在土里腐烂的正是肉身。因为这是天主法令所摧毁的本质：\Quote{你是土，你要归于土}\BibleRef{创 3:19}，因为它取自土。正是出于这一原因，宗徒借用了肉身被\Quote{种下}的说法，因为它归于土，而土是伟大的仓库，供那些要被存放其中并再次从中取出的种子使用。因此，他再次确认了这些经文，盖上（他自己受默感权威的）印记，说道：\Quote{因为经上这样记载}\BibleRef{格前 15:45}，以使你们不可以为\Quote{被种下}意味着别的什么，而不是\Quote{你要回到你所出自的土中}；也不可使\Quote{因为经上这样记载}这句话指涉任何其他事物，而不指肉身。\par

% structure: p0107 source=h2 target=section reason=relative-heading-rank; base=h2

\section{第53章。不是灵魂，而是那已经死了的自然身体，才是要复活的。论及\Person{拉匝禄}的复活。\Person{基督}的复活，作为第二\Person{亚当}，保证我们的复活。}

然而有些人却辩称，灵魂就是\Quote{自然身体}\OriginalTerm{自然身体}{natural (or animate) body}，其目的是要将肉身排除在复活的身体之外。既然那将要复活的身体明确就是那在死亡中所播下的，他们就应当被挑战来审视这事实本身。不然，就请他们证明，灵魂是在死后被播下的；总之，灵魂经历死亡——即被分解、肢解、消散在泥土中，而天主从未对灵魂宣告过任何这类事；请他们向我们展示灵魂的可朽性、羞辱（以及）软弱，这样灵魂也能在不朽、荣耀和权能中复活。\BibleRef{格前 15:42}现在，以拉匝禄为例（我们可将其视为复活的最显著例证），他的肉身软弱地躺卧，他的肉身在\Emph{腐烂}的羞辱中几乎朽坏，他的肉身在腐朽中发臭，然而拉匝禄是作为\Emph{肉身}复活的——无疑带着他的灵魂。但那个灵魂是不朽的；没有人用殓布包裹它；没有人把它放在坟墓里；还没有人闻到它\Quote{发臭}；四天来没有人看见它被\Quote{播下}。那么，这整个状态，这拉匝禄的全部结局，实际上所有人类的肉身仍在经历，但没有任何人的灵魂经历。因此，那宗徒的全部描述显然所指的、他清楚谈论的那个实体，在播下时是自然身体，在复活时是属神身体。为使你这样理解，他指向同一结论，当同样地依据同一段圣经的权威，他给我们显示\Quote{第一个人亚当成了有灵的灵魂}。既然亚当是第一个人，既然也是肉身先于灵魂而成为人，那么无疑，肉身成了有灵的灵魂。而且，既然是一个物质实体采取了这种状态，那么当然就是自然身体成了有灵的灵魂。他们会用什么名称来称呼它呢？除非是它通过灵魂而成为的，除非是它在灵魂之前所不是的，除非是它离开灵魂之后（除非通过复活）永远不能成为的？因为当它重新获得灵魂时，它再次成为自然身体，以便成为属神身体。因为它于复活中只是恢复了它曾一度拥有的状态。因此，灵魂被称为自然身体的理由，完全不像肉身拥有这一称谓那样充足。事实上，肉身在成为自然身体之前，就已经是一个身体。当肉身与灵魂结合时，它就变成了自然身体。现在，虽然灵魂是一个物质实体，但由于它不是有灵的身体，而是赋予生命的身体，所以它不能被称为\OriginalTerm{有灵（或自然）身体}{animate (or natural) body}，也不能成为它所产生的东西。确实，当灵魂附着于其他事物时，它使那事物有灵；但除非这样附着，它如何能产生生命？因此，正如肉身起初因接受灵魂而成为自然身体，同样最终当披上圣神时，它将变成属神身体。现在，宗徒通过在亚当和基督身上分别引出这一顺序，恰当地区分了这两种状态，即它们差异的要点。当他称基督为\Quote{末后的亚当}\BibleRef{格前 15:45}时，你可以由此发现他如何在全部教导中竭力确立肉身的复活，而非灵魂的复活。如此，第一个人亚当是肉身，不是灵魂，只是后来才成为有灵的灵魂；末后的亚当，基督，之所以是亚当，只因为他是人，而作为人，只是作为肉身，而非作为灵魂。因此宗徒接着说：\Quote{然而，属神的不是在先，而是属自然的在先，以后才是属神的}\BibleRef{格前 15:46}，正如在两个亚当的情形中。那么，你不认为他是在同一肉身中区分自然身体和属神身体吗？因为他已经在两个亚当中，即在第一个人和末后的人中，在其中作了区分。因为基督和亚当在哪一个实体上彼此相似？无疑是在他们的肉身上，尽管也可能在灵魂上。然而，是在肉身方面，他们二者都是人；因为肉身先于\Emph{灵魂}而成为人。实际上，正是从肉身，他们能够排列等级，被视为——一个是第一个人，另一个是末后的人，或亚当。此外，性质不同的事物，只有当它们实体的多样性时，才无法被安排在同一顺序中；因为当多样性在于地点、时间或状态时，它们很可能可以一起分类。但在这里，他们被称为首先的和末后的，是基于他们（共同）肉身的实体；正如后来再次，第一个人（被说成）出于地，第二个人出于天；但虽然按圣神祂是\Quote{出于天}，祂却是按肉身成人。既然正是肉身（而非灵魂）使得两个亚当的顺序相容，以致他们之间的区分是\Quote{第一个人成了有灵的灵魂，末后的人成了赋予生命的圣神}，同样，这一区分已经暗示了结论：这一区分是由于肉身；所以，正是关于肉身，这些话说的：\Quote{然而，属神的不是在先，而是属自然的在先，以后才是属神的}。同样，在先前的一段经文中，也必须理解同一\Emph{肉身}：\BibleQuote{所播下的是自然身体，所复活的是属神身体；因为属神的不是在先，而是属自然的在先：既然第一个人亚当成了有灵的灵魂，末后亚当成了赋予生命的圣神。}\BibleRef{格前 15:44}一切都是关于人，\Emph{且}一切都是关于肉身，因为是关于人。\par

那么我们要说什么呢？难道肉身在此生不也因信德而拥有圣神吗？因此仍需追问：如何能说那属魂（或自然）的身体是被播种的？诚然，肉身在此世已经领受了圣神——但\\Emph{仅它的}\\Quote{凭据；}而就灵魂而言，它领受的不是凭据，而是完整的拥有。因此它被称为\\Emph{属魂的}（或自然的）身体，明确地是因为灵魂（或\\OriginalTerm{灵魂}{anima}）这一更高级的实体，它被播种在其中，注定要在将来通过它将要获得的圣神的圆满拥有而成为属神的身体，即复活时所提升的。那么，如果它更常以它所充分具备的实体来命名，而不是以它仅有微末的实体来命名，这有什么可惊奇的？\par

% structure: p0110 source=h2 target=section reason=relative-heading-rank; base=h2

\section{第五十四章 死亡被生命吞灭：这句话在论身体复活中的意义}

再者，疑问往往由零散的术语引发，正如由连贯的句子引发一样。因此，由于宗徒所说：\\Quote{使这必死的被生命吞灭}\\BibleRef{格后 5:4}——指肉身而言——他们曲解\\Emph{吞灭}一词，以为是指肉身的实际消灭；好像我们不能说自己吞下苦胆或吞下忧伤，意思是隐藏和掩饰，并保持在我们内。事实是，当经文写道：\\Quote{这必死的必须穿上不死的}\\BibleRef{格前 15:53}，它解释了\\Quote{必死的被生命吞灭}是什么意思——即，当穿上不死时，它被隐藏、遮盖、包含在其中，而不是被消耗、毁灭、丧失。但你或许会回答我说，照此，死亡在被吞灭后仍然安全。那么，我请你按正确的含义区分形式相似的词。死亡是一回事，必死性是另一回事。死亡被吞灭是一回事，必死性被吞灭是另一回事。死亡不能成为不死的，但必死性可以。此外，既然经上写道：\\Quote{这必死的必须穿上不死的}\\BibleRef{格前 15:53}，它又如何能被生命吞灭呢？但它是如何被生命吞灭的（在被生命毁灭的意义上）？实际上它是被接收、被恢复、被包括在生命之内。至于其他的，死亡被完全毁灭而吞没才是公正合理的，因为死亡本身就怀着同样的意图来吞噬。宗徒说，死亡藉着施展力量吞噬了，因此它自己在争战中被吞噬了，\\Quote{\\Emph{被胜利吞灭}}。\\Quote{死亡啊，你的刺在哪里？死亡啊，你的胜利在哪里？}因此，生命作为死亡的大敌，将在争战中为了救恩而吞灭死亡在争战中所吞噬的毁灭之物。\par

% structure: p0112 source=h2 target=section reason=relative-heading-rank; base=h2

\section{第五十五章 事物状态的改变并非其实体的毁灭。此原则在我们主题上的应用。}

然而，尽管在证明肉体将要复活时，我们\Emph{\OriginalTerm{当然}{ipso facto}}也证明了除了我们所讨论的这肉体之外，没有别的肉体将分享那复活，但孤立的议题及其场合确实需要它们自身的讨论，即使它们已经得到充分的回应。因此，我们将对一种变化的效力与理由作更充分的解释，这种变化如此之大，以致几乎让人以为将要复活的是另一种不同的肉体；仿佛如此巨大的变化就等于完全终止和彻底毁灭先前的自我。然而，必须区分\Emph{变化}（无论多大）与一切具有\Emph{毁灭}性质的事物。因为发生变化是一回事，被毁灭是另一回事。如果肉体遭受的变化达到毁灭的地步，这种区分便不再存在。然而，若不以复活中将显现的转变状态下自身持续存在，它必被变化所毁灭。因为正如不复活便消亡，即便复活，若在变化中丧失，也同样消亡。它将完全失去未来的存在，如同根本没有复活。为了不再存在而复活，这何其荒谬，尽管它本可以不复活而失去存在——因为它已经开始不存在了！绝对不同的事物，如变化与毁灭，不容混合与混淆；它们的作用也不同：一者毁灭，一者改变。因此，被毁灭的并未改变，改变的并未毁灭。消亡是完全停止曾是的东西，而改变是以另一种状态存在。若某物以另一种状态存在，它仍然是自身；因为它未消亡，便依然存在。它经历了变化，而非毁灭。某物可经历彻底变化，却仍保持为同一物。同样，人在现世生活中，其本质也完全可以是自身，同时经历各种变化——习惯、体型、健康、境遇、尊严、年龄、趣味、事业、财产、法律和习俗——但丝毫不失人性，也不会变成另一个人而不再相同；其实我几乎不该说\Quote{另一个人}，而是\Quote{另一物}。圣经也给我们提供了这种变化的实例。梅瑟的手变了，变得如死一般，无血无色，冰冷僵硬；但恢复温暖和天然颜色后，它又成了同样的血肉。\BibleRef{出 4:6}此后，同一位梅瑟的脸庞变了，发出目不能视的光辉。但他仍是梅瑟，即使不可见。同样，斯德望已显出天使的面容，\BibleRef{宗 6:15}尽管跪在石击之下的是他的人腿。\BibleRef{宗 7:59}主又在登山独处时，将自己的衣袍变成光明之袍，但仍保留伯多禄能认出的面貌。\BibleRef{玛 17:2}在那同一场景中，梅瑟和厄里亚也证明了同样的身体状态可延续于光荣中——一人尚未得回肉身的样式，另一人尚未脱去肉身的真实。正是充满这一光辉的榜样，保禄曾说：\BibleQuote{祂将改变我们卑贱的身体，使之相似祂光荣的身体。}\BibleRef{斐 3:21}但若你坚持认为变形和转变等于任何本质的消灭，那么随之而来的就是，\BibleQuote{撒乌耳变成了另一个人}\BibleRef{撒上 10:6}时，便脱离了他自身的身体本质；并且撒殚自己\BibleQuote{装作光明的天使}\BibleRef{格后 11:14}时，就失去了其自身的特性。这并不是我的意见。同样，变化、转变和改造必然会为了带来复活而发生，但\Emph{肉体的}本质仍将安全地保存。\par

% structure: p0114 source=h2 target=section reason=relative-heading-rank; base=h2

\section{第五十六章 最后审判的过程及其赏罚，唯有在复活的肉身与我们现世肉身同一的基础上才有可能。}

因为，如果让一个实体做工，另一个实体得酬报——即我们的这肉身因殉道而破碎，另一肉身戴冠冕；或者反过来，我们的这肉身在污秽中打滚，另一肉身受惩罚——这是多么荒谬，实则多么不义，而在这两方面又多么不相称于天主！与其轻慢天主的智慧与公义，不如立即全然放弃对复活的希望；\Person{马西昂}复活强于\Person{瓦伦提尼安}复活，因为不能相信现世之人的心神、记忆或良心会被披上不朽与不坏之改变的外衣所消除；否则，复活的一切得益与果实，以及天主对灵魂和身体审判的持久效果，必定落空。若我记不起是我侍奉了祂，我怎能将光荣归于天主？我若不知是我欠祂感谢，怎能向祂歌唱\Quote{新歌}？为什么要唯独对肉体的变化提出异议，而不是也对灵魂的变化提出，灵魂在所有方面都比肉体优越？为何发生这种情况，即那在现世肉体中经历人生全程、学习认识天主、穿上基督、在此肉体中播下救恩希望的同一位灵魂，却要在我们所不知的另一肉体中收获？那肉体必定是极其蒙恩，能如此白白地享受生命。但若灵魂也不改变，那么就没有灵魂的复活；除非它成为某种不同之物，否则不会被相信已复活。\par

% structure: p0116 source=h2 target=section reason=relative-heading-rank; base=h2

\section{第五十七章 我们的身体，无论在死前或死后如何伤残，在复活时将恢复完全的完整。获释奴隶的例证。}

我们现在来到不信者最通常的挑剔。他们说，如果真是同一实体带着它的一切形状、线条和品质被召回\Emph{生命}，那么为什么不是带着它的一切其他特征呢？那么瞎子、瘸子、瘫子以及任何带着显著记号去世的人，都将带着同样的标记归来。现在事实是什么呢？尽管你因自大而藐视从天主接受如此浩大的恩宠？难道不是这样：当你现在只承认灵魂的救恩时，你赋予人的救恩是以他们一半的本性为代价吗？相信复活有什么益处，除非你的信德涵盖它的全部？如果肉体在分解之后还要被修复，那么在一次猛烈伤害之后就更会被复原了。更大的情况为较小的定下规则。肢体的截除或压碎不就是那肢体的死亡吗？那么，如果整个人的死亡因复活而被撤销，我们对他身体一部分的死亡该说什么呢？如果我们为光荣而改变，那么为完整而改变岂不更多！我们身体所遭受的任何损失都是偶发的，而它们的完整是其自然属性。我们是在这种状态下出生的。即使我们在母胎中受伤，这也是一个已经存在的人所遭受的损失。自然状态先于伤害。正如生命由天主赐予，同样也由祂恢复。我们接受它时是什么样，我们恢复它时就是什么样。我们被恢复到自然状态，而非伤害状态；我们复活到我们出生时的状态，而非意外造成的状态。如果天主不使整个人复活，那么祂就没有复活死人。因为哪个死人能说是完整的，即使他是完整地死去？谁没有受伤，那不就是没有生命吗？哪具身体在死亡时、冰冷时、骇人时、僵硬时、成为尸体时是无损伤的？一个人什么时候比他完全虚弱时更虚弱？什么时候比他完全不能动时更瘫痪？因此，死人复活完全不亚于他被恢复到完整的状态——以免他，确实，在他还未复活的那一部分中仍然死去。天主完全能够重新制造祂曾制造的东西。祂的这种能力和这种毫不吝啬的恩宠，已经在基督里充分保证了；并且向我们（在祂内）显示了自己不仅是肉体的恢复者，也是其裂痕的修补者。因此宗徒说：\BibleQuote{死人必要复活起来，成为不朽坏的}（或不受损伤的）。\BibleRef{格前 15:52} 但怎么会这样呢？除非他们成为完整的，这些人要么因健康丧失而消瘦，要么因坟墓的长期腐朽而耗损？因为他提出两个条款：\BibleQuote{这可朽坏的必须穿上不朽坏，这可死的必须穿上不死}，\BibleRef{格前 15:53} 他不是重复同一句话，而是指明一个区别。因为，他把\Emph{不死}归于不再有死亡，把\Emph{不朽坏}归于耗损身体的修复，他使一个适合于复活，另一个适合于身体\Emph{的恢复}。再者，我想他许诺给得撒洛尼人整个人性的完整。所以对于伟大的未来，无需害怕残缺或有缺陷的身体。完整，无论是保守或恢复的结果，在它收回它所失去的一切之后，将不能再失去什么。现在，当你争辩说，如果同样的肉体必须复活，那么肉体仍将遭受同样的痛苦时，你鲁莽地将本性置于她的主之对立面，且不敬地把她\Emph{的}律法与祂\Emph{的}恩宠相对比；好像上主天主不被允许既改变本性，又保存她，而不受律法的约束。那么，我们如何读到：\BibleQuote{在人不可能，在天主却不然；} \BibleRef{玛 19:26} 以及：\BibleQuote{天主拣选了世上愚拙的，叫有智慧的蒙羞？} \BibleRef{格前 1:27} 我问你，如果你释放你的奴隶（既然同样的肉体和灵魂仍留给他，这些曾暴露于鞭子、脚镣和鞭打之下），那么他适合再受同样的旧痛苦吗？我想不。他反而得到白袍的恩宠、金戒指的厚爱以及其恩主的名字、族系和餐桌之荣耀。那么，将同样的特权给予天主，藉着这样的改变，改造我们的状况，而非我们的本性，即夺去它的一切痛苦，并用保护的保障围绕它。因此，我们的肉体即使在复活后仍将存留——它确实是可感受痛苦的，因为它是肉体，而且也是同一肉体；但同时又是不可受苦的，因为它已被主释放，正是为了这个目的和意图，即不再能够忍受痛苦。\par

% structure: p0118 source=h2 target=section reason=relative-heading-rank; base=h2

\section{第五十八章 从我们复原身体的这种完美中，将涌出不被打扰的喜乐与平安的意识。}

依撒意亚说：\Quote{永远的喜乐}要临到他们头上。\BibleRef{依 35:10} 嗯，在复活之前，没有什么永恒。他接着说：\Quote{悲哀和叹息}要逃遁。天使也向若望重复同样的话：\Quote{天主要擦去他们眼中一切的眼泪；}\BibleRef{默 7:17} 就是那曾经流泪的同一双眼睛，它们本可能再次流泪，如果天主的慈爱不干涸一切泪泉的话。又说：\Quote{天主要擦去他们眼中一切的眼泪；不再有死亡；}\BibleRef{默 21:4} 因此也不再有任何腐朽，因为不朽赶走了腐朽，正如不死被永生所取代。如果悲哀、忧伤、叹息和死亡本身因灵魂和肉身的痛苦而攻击我们，那么除了这些原因的停止，即肉身和灵魂的痛苦，它们怎能被除去呢？在天主面前，你会在哪里找到逆境？在基督的怀抱中，哪里会有敌人的入侵？在圣神的面前，哪里会有魔鬼的攻击？——现在魔鬼和他的使者已被\Quote{扔在火湖里}。现在哪里还有必要，以及他们所谓的命运或定数？在永恒赦免之后，还有什么灾祸等待着从死亡中得救的人？在恩宠之后，还有什么愤怒对待被和好的人？在重新得力之后，还有什么软弱？在获得救恩之后，还有什么危险和恐惧？以色列子民的衣服和鞋子在四十年间没有磨损，依然如新；\BibleRef{申 29:5} 在他们身上，合宜和适当的点抑制了指甲和头发的过度生长，这样任何过分之处都不会被归于不雅；巴比伦的火没有烧伤三位青年的头巾或裤子，尽管这种服饰对犹太人来说多么异样；\BibleRef{达 3:27} 约纳被深渊的怪兽吞下，它的肚子能吞下整艘船，三天后又被完好地吐出来；哈诺客和厄里亚，他们至今没有经历复活（因为他们甚至没有遇到死亡），正在充分学习肉身免于一切屈辱、一切损失、一切伤害和一切耻辱意味着什么——他们被从这个世界提升，并因此已经成为永生的候补者；——这些显著的事实为哪种信仰作证，如果不是那种应当在我们心中激发信念，认为它们是我们未来完善\Emph{和完美复活}的证明和凭证呢？因为，借用宗徒的话，这些是\Quote{我们的预像}；\BibleRef{格前 10:6} 并且它们被记载下来，好让我们相信：主比一切关于身体的自然律更强大，并且祂更强调地表明自己是肉身的保存者，因为祂甚至为肉身保存了衣服和鞋子。\par

% structure: p0120 source=h2 target=section reason=relative-heading-rank; base=h2

\section{第五十九章 我们的肉身在复活时，能够在不失去其本质特性的情况下，承受永生或永死的改变条件。}

但是，你反驳说：将来的世界带有一种不同安排的特性，甚至是永恒的；因此，你主张，今生的非永恒的本质不可能拥有如此不同特征的状态。这固然是真的，如果人是为将来的安排而造，而不是安排为人而造的话。然而，宗徒在他的书信中说：\Quote{无论是世界，或是生命，或是死亡，或是现在的事，或是将来的事，一切都是你们的；}\BibleRef{格前 3:22} 他在这里使我们成为将来世界的继承者。依撒意亚帮不了你，当他说：\Quote{凡有血气的，尽都如草；}\BibleRef{依 40:7} 而在另一处又说：\Quote{凡有血气的，都要看见天主的救恩。} 他所区分的是人的结局，而不是人的本质。但谁不认为天主的审判在于双重判决，即救恩和惩罚呢？因此，\Quote{凡有血气的，尽都如草}，注定被丢进火里；而\Quote{凡有血气的，都要看见天主的救恩}，被预定得永生。就我自己而言，我相当确信，我是在我自己的肉身中犯了奸淫，也正是在我自己的肉身中我努力追求节制。如果有人在自己身上携带着两件好色的器具，他当然有权割去那不洁肉身的\Quote{草}，只为自己保留那将要看见天主治恩的肉身。但是，当同一位先知向我们描述万民有时被视为\Quote{天平上的微尘}\BibleRef{依 40:15}，以及\Quote{虚无和空虚}，有时又将盼望和\Quote{仰望主的名}和祂的膀臂，我们是否因着\Emph{言辞的多样性}而对外邦民族产生误解？难道有些民族因本质不同而信主，另一些则被视为尘土？相反，基督作为真光，照耀了海洋界限内的万民，并从我们头顶的诸天之上照耀。为什么，甚至在这地上，华伦提诺派就已经学习了他们的谬误；在信主的民族和不信的民族之间，就其身体和灵魂而言，不会有任何状态的不同。因此，正如祂在同一个民族之间作出了状态上的区分，而非本质上的区分，同样，祂也区分了他们的肉身——这肉身在这些民族中是同一本质——不是按照他们的物质结构，而是按照他们功过的报酬。\par

% structure: p0122 source=h2 target=section reason=relative-heading-rank; base=h2

\section{第六十章 我们身体的一切特征——性别、各种肢体等——都将被保留，无论其功能可能发生什么变化，然而关于这一点我们无法判断。以修复的船为类比。}

但请看，他们多么执拗地不断堆积对肉身的攻讦，尤其是针对其同一性，甚至从我们肢体的功能中引出论据；一方面说，这些肢体应当永久地继续它们的劳作和享乐，作为同一身体框架的附属物；另一方面又争辩说，既然肢体的功能终有一日会停止，身体框架本身就必须被摧毁，因为没有了肢体，它的持续存在被认为像没有了功能的肢体本身一样不可想象。他们问，那时口腔的腔体、成排的牙齿、咽喉的通道、胃的分支、肚腹的深渊、以及缠绕的肠组织还有什么用处，既然再也没有进食和饮水的余地？这些器官还有什么要摄取、咀嚼、吞咽、分泌、消化、排泄？当对食物的一切挂虑都停止时，我们的双手、双脚以及所有劳作的肢体又有何益？当同居、怀孕和哺育婴儿都停止时，意识到精液分泌的腰胯、两性的所有其他生殖器官、胚胎的工场、以及乳房的源泉又能起什么作用？简言之，当整个身体变得无用之时，整个身体还有什么用处？作为对这一切的回应，我们此前已经确立了原则：未来状态的分施不应与现世相比较，并且在两者之间会发生改变；现在我们再补充一点：我们身体肢体的这些功能将继续供给现世生命的需要，直到生命本身从时间过渡到永恒，正如属灵的身体取代属血气的身体，直到 \BibleQuote{这个可朽坏的必须穿上不朽坏的，这个可腐烂的必须穿上不腐烂的：} \BibleRef{格前 15:53}；因此，当生命本身摆脱一切需要时，我们的肢体也将摆脱它们的服务，因而不再被需要。不过，虽然从职责中解脱出来，它们仍将被保存以待审判，\BibleQuote{好使每一个人都照他身体所行的受报应。} \BibleRef{格后 5:10} 因为天主的审判座要求人保持完整。然而，没有肢体他无法保持完整，因为他是由肢体的实质构成的，而非由功能构成；除非你胆敢断言，一艘船没有龙骨、船首、船尾，没有整个船身的坚固，仍是完美的。然而，我们多少次看到同一艘船，经历了风暴的摧残和腐朽的侵蚀，修理并恢复了所有木材，以崭新的船体优美地破浪航行！难道我们还要怀疑天主的技艺、意愿和权力吗？此外，如果一位富有的船主，不惜金钱只是为了娱乐或展示，彻底修好了他的船，然后选择不再让它航行，你会争辩说，尽管船只不再用于实际服务，但旧的形式和完工对船仍不必要吗？因为即使不考虑服务，单是船只的安全就需要这样的完整性。因此，我们在此足够了考虑的唯一问题是：主在安排人的救恩时，是否定意为他的肉身预备救恩？祂是否愿意同一个肉身被更新？\Emph{如果如此}，那么你以未来状态中肢体的无用性来判定肉身不能更新，就是不恰当的。因为一件事物可以被更新，但可能因\Emph{无事可做}而仍然无用；但不能说它无用，如果它根本就不存在。如果它确实存在，它完全有可能不是无用的；它\Emph{也许会有事可做}；因为在天主面前不会有懒惰。\par

% structure: p0124 source=h2 target=section reason=relative-heading-rank; base=h2

\section{第六十一章 我们身体性别的细节，以及我们各种肢体的功能。为异端强迫搜罗所有无耻狡辩的必要性辩护。}

如今，人啊，你接受你的口，是为了吃你的食物，喝你的饮料；为何不为了更高的目的，即发出言语，使你与所有其他动物分别出来呢？为何不更为了宣讲\Emph{天主的福音}，使你甚至可以成为祂的司铎，在众人面前作祂的辩护人呢？亚当在摘树上的果子之前，曾给动物各起名字；在吃之前，他就预言了。再者，你接受你的牙齿，是为了吃东西；为何不更为了在每次开口，无论大小，以适当的防御来装饰你的嘴呢？为何不也为了节制你舌头的冲动，保护你清晰的言语免于失败和粗暴呢？让我告诉你，（如果你不知道），世上有没牙齿的人。看看他们，问问是否连一个牙齿的笼子也不是嘴的荣耀。男人和女人的下体有出口，他们无疑通过它们满足肉欲；但它们为何不被视为清洁排放体液的自然通道呢？此外，女人体内有容器容纳人的种子；但它们不是为分泌那些血液的流出物而设计的吗？因为她们较迟缓、较软弱的性别不足以排出这些。就连这些细节也必须提及，因为\Emph{\OriginalTerm{异端者}{heretics}}挑出我们身体中适合他们的部分，轻率地处理它们，并随其奇想，对我们肢体的自然功能倾泻嘲弄和蔑视，目的是推翻复活，使我们因他们的诡辩而羞愧；他们没有考虑到，在功能停止之前，它们的原因本身就会消逝。将不再有食物，因为不再有饥饿；不再有饮料，因为不再有干渴；不再有同居，因为不再有生育；不再有吃喝，因为不再有劳苦和劳累。死亡也将停止；因此不再需要食物的营养来维持生命，也不再有\Emph{母亲的}肢体为延续人类而承受负担。但即使在现世生活中，我们的胃和生殖器官也可能有停止其职能的时候。因为梅瑟曾四十天\BibleRef{出 24:8}和厄里亚\BibleRef{列上 19:8}禁食，只靠天主生活。因为那原则早已被祝圣：\Quote{人生活不只靠饼，也靠天主口中所发的一切言语。}看这里我们未来力量的模糊轮廓！我们甚至尽可能地使我们的口不吃食物，使我们的性器官脱离结合。有多少自愿守贞的人！有多少与基督订亲的贞女！有多少男女天生不育，结构不能生育！如果即使在地上，我们肢体的功能和快乐都可能暂停，这种中断像安排本身一样只能是暂时的，而人的安全仍然不受损害，那么当他的救恩稳固，尤其是在永恒的安排中，我们岂不更将停止渴望那些甚至在地上我们也不习惯克制其欲求的事物吗？\par

% structure: p0126 source=h2 target=section reason=relative-heading-rank; base=h2

\section{第六十二章 我们在复活的光荣生命中预定肖似天使。}

然而，我们的主的一句话为这个讨论画上了句号：祂说：\Quote{他们要像天使一样。}正如因为不死而不结婚，同样也因为不必屈服于我们身体状态的任何类似需要；甚至天使有时也\Quote{像}人一样，通过吃喝，并让他们的脚接受洗脚——他们披上人的外表，却没有失去自己的本性。因此，如果天使当成为人时，在自己不变的灵性实体中，如同肉身一样被对待，那么人当成为\Quote{像天使一样}时，为什么不能在自己不变的肉身实体中，经历灵性存在的对待，不再受肉体的惯常诱惑（当他们穿上天使的衣袍时），正如天使曾（披上人形时）不再受灵性的诱惑一样呢？因此，我们不会因为不再受肉体的惯常需要困扰而停止存留在肉身中；正如天使不会因为其灵性属性的暂停而停止存留在他们的灵性实体中一样。最后，基督没有说：\Quote{他们要成为天使，}以免废除他们作为人的存在；而是说：\Quote{他们要像天使一样，}以便保持他们的人性不受损害。当祂将天使的相似归于肉身时，并没有剥夺其固有的实体。\par

% structure: p0128 source=h2 target=omitted reason=duplicate-page-title

因此，肉身必要复活，完全在每个人内，以它自身的身份，以它绝对的完整。无论它在何处，都在天主面前的安全保管中，借着那最忠信的\Quote{介于天主与人之间的\OriginalTerm{中保}{Mediator}，就是那人基督耶稣，} \BibleRef{弟前 2:5} 祂将使天主与人、人与天主和好，并使灵与肉身、肉身与灵和好。祂已经在自己内将两种本性联合；祂将它们配合起来，如同新娘和新郎，在婚姻生活的相互纽带中。现在，如果有人坚持将灵魂当作新娘，那么肉身将作为她的嫁妆跟随灵魂。灵魂绝不可成为被遗弃者，被新郎赤身裸体地带回家。她在肉身中有她的嫁妆、她的装扮、她的财富，这些将像乳姐妹的爱与忠诚一样陪伴她。但假设肉身是新娘，那么在基督耶稣里，她在祂血的盟约中，收到了祂的圣神作为她的配偶。现在，你以为她的消失，你可以确定只是暂时的退隐。不仅灵魂从视野中消失。肉身也有她暂时的离去——在水中、在火中、在飞鸟中、在野兽中；她似乎被溶解在这些之中，但她只是被倾倒进它们之中，如同倒入容器。如果这些容器后来也无法容纳她，她甚至从这些中逃脱，回到她的母亲大地，她仿佛再次被大地的隐秘拥抱所吸收，最终将显现在眼前，如同亚当被召来听他的主和造物主的话时：\Quote{看，人已经成了我们中的一个！} \BibleRef{创 3:22}——那时她彻底\Quote{认识}了\Quote{她所逃脱的恶}，和\Quote{她所获得的善}。那么，灵魂啊，你为什么要嫉妒肉身呢？除了主以外，没有谁是你应该如此深爱的；没有谁比你更像兄弟，它甚至是与你自己一同在天主内诞生的。你更应该通过你的祈祷为她求得复活：她的罪，无论是什么，都是因你而来。然而，如果你恨她，并不奇怪；因为你已经否认了她的创造者。你已经习惯于否认或改变她在基督内的存在——要么通过删减或曲解圣经，从而败坏那降生成人的天主圣言本身，并且，尤其是引入伪经的奥秘\Emph{和}亵渎的寓言。然而，全能天主，以祂最恩慈的天主眷顾，借着\Quote{在末日将祂的圣神倾注在一切血肉之上，在祂的仆人和婢女身上，}遏制了这些不信和悖逆的骗术，重振了人们对肉身复活的动摇信德，并以他们神圣言语和意义的清晰之光，扫清了古代圣经（天主的两个盟约）的一切晦涩和歧义。现在，既然\Quote{异端是必然要有的，好使那些被认可的人显明出来；} \BibleRef{格前 11:19} 然而，既然这些异端若没有圣经的些许支持就无法装出大胆的姿态，因此十分明显，古代圣经为它们邪恶的教义提供了各种材料，而正是这些材料（如此扭曲）可以从同一圣经中被驳斥。因此，适得其所，圣神不应再保留祂恩宠之光的倾注于这些默感著作上，以便它们能够传播\Emph{真理的}种子，没有混杂异端的诡辩，并从其中拔除它们的莠子。因此，祂现在已经驱散了过去的全部困惑，和他们自行选择的寓意和比喻，通过对全部奥秘的公开而清晰的解释，借着从护慰者涌出的丰富水流般的新预言。只要你从祂的泉源中取水，你就永不会渴慕其他教义：对微妙问题的狂热渴望不会再消耗你；\Emph{而是}通过不断饮用肉身的复活，你将因这清凉的饮品而得满足。\par

\SourceRecord{Translated by Peter Holmes.；Ante-Nicene Fathers；Vol. 3.；Edited by Alexander Roberts, James Donaldson, and A. Cleveland Coxe.；Buffalo, NY: Christian Literature Publishing Co.,；1885.；<http://www.newadvent.org/fathers/0316.htm>.}
